\section{AIME II 2011}
        \begin{enumerate}
        	\item (\textit{AIME II 2011}) Gary purchased a large beverage, but only drank $m/n$ of it, where $m$ and $n$ are relatively prime positive integers. If he had purchased half as much and drunk twice as much, he would have wasted only $2/9$ as much beverage. Find $m+n$.


 
 \textbf{Answer: }037



\textbf{Solution}

Let $x$ be the \href{https://artofproblemsolving.com/wiki/index.php?title=Fraction}{fraction} consumed, then $(1-x)$ is the fraction wasted. We have $\frac{1}{2} - 2x =\frac{2}{9} (1-x)$, or $9 - 36x = 4 - 4x$, or $32x = 5$ or $x = \frac{5}{32}$. Therefore, $m + n = 5 + 32 = \boxed{037}$.

 





\textbf{Solution 2}


 WLOG, Gary purchased \( n \) liters and consumed \( m \) liters.   After this, he purchased \( \frac{n}{2} \) liters, and consumed \( 2m \) liters.   He originally wasted \( n-m \) liters, but now he wasted \( \frac{n}{2} - 2m \).   \[ \frac{n}{2} - 2m = \frac{4}{18} \cdot (n-m) \] \[ 9n - 36m = 4n - 4m \implies 5n = 32m \implies \frac{m}{n} = \frac{5}{32}. \]

 Thus, the answer is $\boxed{37}$~idk123456

 





\textbf{Solution 3}

Say Gary buys 2 litres, drinks $x$ of that. So, he wastes $2-2x$ of the total solution. Now, say he buys 1 litre, drinks $2 \cdot 2x$ that. Then, he only wastes $\frac{2}{9}$ that. Notice we do $2 \cdot 2x$ and not $2x$, as he doesn't drink $x$, rather, he drinks $x$ THE amount of beverage, which is $2x$, so drinking double that of before is drinking double that of $2x$, or $4x$. So, $1-4x = \frac{2}{9}(2-2x)$. Solving, $x  = \frac{5}{32} \implies \boxed{37}$.

 




	\item (\textit{AIME II 2011}) On square $ABCD$, point $E$ lies on side $AD$ and point $F$ lies on side $BC$, so that $BE=EF=FD=30$. Find the area of the square $ABCD$.


 
 \textbf{Answer: }810



\textbf{Solution}

Drawing the square and examining the given lengths,
                    \begin{figure}[H]
                        \centering
                        \includegraphics[width=0.5\textwidth]{Solution/AIME_II_2011_2_Solution_Figure_1.png}
                    \end{figure}
                    you find that the three segments cut the square into three equal horizontal sections. Therefore, ($x$ being the side length), $\sqrt{x^2+(x/3)^2}=30$, or $x^2+(x/3)^2=900$. Multiplying both sides by $9$ and simplifying, we find that $10x^2=8100$. Dividing by ten gives $x^2=810.$

 Area of the square is $\fbox{810}$.

 





\textbf{Solution 2}

Drop horizontal $E$ down to $BC$, $F$ down to $AD$, called lines $l, m$ respectively. $BE$ is a reflection upon $l$, which is a reflection upon $M$. Define $l$ to intersect $BC$ at $G$, and apply with $m$ and $AD$ at $H$. Understand that $ABGE, GEFH, HFCD$ are all with diagonal 30. Then, the three rectangles are simply reflected over their side lengths, which then form the area of the square, so each of the rectangles share $\frac{1}{3}$ the area of the square. Due to their trisections, their sides also trisect the square (exception of the longer side), so let that be $a$. Then, $10a^2 = 900$, so $a = \sqrt{90}$. We require $9a^2 = 9 \cdot 90 = \boxed{810}$.

 


 




	\item (\textit{AIME II 2011}) The degree measures of the angles in a convex 18-sided polygon form an increasing arithmetic sequence with integer values. Find the degree measure of the smallest angle.


 
 \textbf{Answer: }143



\textbf{Solution 1}

The average angle in an 18-gon is $160^\circ$. In an arithmetic sequence the average is the same as the median, so the middle two terms of the sequence average to $160^\circ$. Thus for some positive (the sequence is increasing and thus non-constant) integer $d$, the middle two terms are $(160-d)^\circ$ and $(160+d)^\circ$. Since the step is $2d$ the last term of the sequence is $(160 + 17d)^\circ$, which must be less than $180^\circ$, since the polygon is convex. This gives $17d < 20$, so the only suitable positive integer $d$ is 1. The first term is then $(160-17)^\circ = \fbox{143}.$

 





\textbf{Solution 2}

Another way to solve this problem would be to use exterior angles. \href{https://artofproblemsolving.com/wiki/index.php?title=Exterior\_angle}{Exterior angles} of any polygon add up to $360^{\circ}$. Since there are $18$ exterior angles in an 18-gon, the average measure of an exterior angles is $\frac{360}{18}=20^\circ$. We know from the problem that since the exterior angles must be in an arithmetic sequence, the median and average of them is $20$. Since there are even number of exterior angles, the middle two must be $19^\circ$ and $21^\circ$, and the difference between terms must be $2$. Check to make sure the smallest exterior angle is greater than $0$: $19-2(8)=19-16=3^\circ$. It is, so the greatest exterior angle is $21+2(8)=21+16=37^\circ$ and the smallest interior angle is $180-37=\boxed{143}$.

 





\textbf{Solution 3}

The sum of the angles in a 18-gon is $(18-2) \cdot 180^\circ = 2880 ^\circ.$ Because the angles are in an arithmetic sequence, we can also write the sum of the angles as $a+(a+d)+(a+2d)+\dots+(a+17d)=18a+153d,$ where $a$ is the smallest angle and $d$ is the common difference. Since these two are equal, we know that $18a+153d = 2880 ^\circ,$ or $2a+17d = 320^\circ.$ The smallest value of $d$ that satisfies this is $d=2,$ so $a=143.$ Other values of $d$ and $a$ satisfy that equation, but if we tried any of them the last angle would be greater than $180,$ so the only value of $a$ that works is $a=\boxed{143}$.

 Note: The equation $2a+17d = 320^\circ$ can also be obtained by using the sum of an arithmetic sequence formula $\frac{2a_1+(n-1)d}{2} \cdot n$. We set $n = 18$ and equate it to 2880, thereby achieving the same result.~Eclipse471~note by cxsmi

 





\textbf{Solution 4}

Each individual angle in a $18$-gon is $\frac {(18-2) \cdot 180^\circ}{18} = 160^\circ$. Since no angle in a convex polygon can be larger than $180^\circ$, the smallest angle possible is in the set $159, 161, 157, 163, 155, 165, 153, 167, 151, 169, 149, 171, 147, 173, 145, 175, 143, 177$. 

 Our smallest possible angle is $\boxed {143}$

 





\textbf{Solution 5 (Variation of [3])}

The sum of all angles in an octadecagon is 2880. Like before, let $a$ be the smallest measure and $d$ be the common difference, so $18a+153d = 2880$. In other words, $a + 8.5d = 160$ so $2a+17d = 320$. Recall $\gcd(2,17) = 1$ which divides and therefore we have solutions. Note that $d = 0$ and $a = 160$ works, so all solutions take the form $a = 160-17k, d = 2k$. Then, the smallest $k$ is $k = 2$, in which $a = 160-34 = \boxed{143}$. No other solutions are valid as maximal $a$ is $a_{\text{max}} = 160-17k + 17d$ in which $109 + 102 > 180$ breaking convexity.

 




	\item (\textit{AIME II 2011}) In triangle $ABC$, $AB=20$ and $AC=11$. The angle bisector of angle $A$ intersects $BC$ at point $D$, and point $M$ is the midpoint of $AD$. Let $P$ be the point of intersection of $AC$ and the line $BM$. The ratio of $CP$ to $PA$ can be expressed in the form $\frac{m}{n}$, where $m$ and $n$ are relatively prime positive integers. Find $m+n$. 


 
 \textbf{Answer: }051



\textbf{Solution 1}


                    \begin{figure}[H]
                        \centering
                        \includegraphics[width=0.5\textwidth]{Solution/AIME_II_2011_4_Solution_Figure_1.png}
                    \end{figure}
                     Let $D'$ be on $\overline{AC}$ such that $BP \parallel DD'$. It follows that $\triangle BPC \sim \triangle DD'C$, so 
 \[\frac{PC}{D'C} = 1 + \frac{BD}{DC} = 1 + \frac{AB}{AC} = \frac{31}{11}\] by the \href{https://artofproblemsolving.com/wiki/index.php?title=Angle\_Bisector\_Theorem}{Angle Bisector Theorem}. Similarly, we see by the Midline Theorem that $AP = PD'$. Thus, 
 \[\frac{CP}{PA} = \frac{1}{\frac{PD'}{PC}} = \frac{1}{1 - \frac{D'C}{PC}} = \frac{31}{20},\] and $m+n = \boxed{51}$.

 





\textbf{Solution 2 (mass points)}

Assign \href{https://artofproblemsolving.com/wiki/index.php?title=Mass\_points}{mass points} as follows: by Angle-Bisector Theorem, $BD / DC = 20/11$, so we assign $m(B) = 11, m(C) = 20, m(D) = 31$. Since $AM = MD$, then $m(A) = 31$, and $\frac{CP}{PA} = \frac{m(A) }{ m(C)} = \frac{31}{20}$, so $m+n = \boxed{51}$.

 





\textbf{Solution 3}

By \href{https://artofproblemsolving.com/wiki/index.php?title=Menelaus\%27\_Theorem}{Menelaus' Theorem} on $\triangle ACD$ with \href{https://artofproblemsolving.com/wiki/index.php?title=Transversal}{transversal} $PB$, 
 \[1 = \frac{CP}{PA} \cdot \frac{AM}{MD} \cdot \frac{DB}{CB} = \frac{CP}{PA} \cdot \left(\frac{1}{1+\frac{AC}{AB}}\right) \quad \Longrightarrow \quad \frac{CP}{PA} = \frac{31}{20}.\] So $m+n = \boxed{051}$.

 





\textbf{Solution 4}

We will use barycentric coordinates. Let $A = (1, 0, 0)$, $B = (0, 1, 0)$, $C = (0, 0, 1)$. By the \href{https://artofproblemsolving.com/wiki/index.php?title=Angle\_Bisector\_Theorem}{Angle Bisector Theorem}, $D = [0:11:20] = \left(0, \frac{11}{31}, \frac{20}{31}\right)$. Since $M$ is the midpoint of $AD$, $M = \frac{A + D}{2} = \left(\frac{1}{2}, \frac{11}{62}, \frac{10}{31}\right)$. Therefore, the equation for line BM is $20x = 31z$. Let $P = (x, 0, 1-x)$. Using the equation for $BM$, we get 
 \[20x = 31(1-x)\]
 \[x = \frac{31}{51}\] Therefore, $\frac{CP}{PA} = \frac{1-x}{x} = \frac{31}{20}$ so the answer is $\boxed{051}$.

 





\textbf{Solution 5}

Let $DC=x$. Then by the Angle Bisector Theorem, $BD=\frac{20}{11}x$. By the Ratio Lemma, we have that $\frac{PC}{AP}=\frac{\frac{31}{11}x\sin\angle PBC}{20\sin\angle ABP}.$ Notice that $[\triangle BAM]=[\triangle BMD]$ since their bases have the same length and they share a height. By the sin area formula, we have that 
 \[\frac{1}{2}\cdot20\cdot BM\cdot \sin\angle ABP=\frac{1}{2}\cdot \frac{20}{11}x\cdot BM\cdot\sin\angle PBC.\] Simplifying, we get that $\frac{\sin\angle PBC}{\sin\angle ABP}=\frac{11}{x}.$ Plugging this into what we got from the Ratio Lemma, we have that $\frac{PC}{AP}=\frac{31}{20}\implies\boxed{051.}$

 





\textbf{Solution 6 (quick Menelaus)}

First, we will find $\frac{MP}{BP}$. By Menelaus on $\triangle BDM$ and the line $AC$, we have 
 \[\frac{BC}{CD}\cdot\frac{DA}{AM}\cdot\frac{MP}{PB}=1\implies \frac{62MP}{11BP}=1\implies \frac{MP}{BP}=\frac{11}{62}.\]This implies that $\frac{MB}{BP}=1-\frac{MP}{BP}=\frac{51}{62}$. Then, by Menelaus on $\triangle AMP$ and line $BC$, we have
 \[\frac{AD}{DM}\cdot\frac{MB}{BP}\cdot\frac{PC}{CA}=1\implies \frac{PC}{CA}=\frac{31}{51}.\]Therefore, $\frac{PC}{AP}=\frac{31}{51-31}=\frac{31}{20}.$ The answer is $\boxed{051}$. -brainiacmaniac31

 





\textbf{Solution 7 (Visual)}


                    \begin{figure}[H]
                        \centering
                        \includegraphics[width=0.5\textwidth]{Solution/AIME_II_2011_4_Solution_Figure_2.png}
                    \end{figure}
                    \textbf{vladimir.shelomovskii@gmail.com, vvsss}

 


 





\textbf{Solution 8 (Cheese)}

Assume $ABC$ is a right triangle at $A$. Line $AD = x$ and $BC = \frac{-11}{20}x + 11$. These two lines intersect at $D$ which have coordinates $(\frac{220}{31},\frac{220}{31})$ and thus $M$ has coordinates $(\frac{110}{31},\frac{110}{31})$. Thus, the line $BM = \frac{11}{51} \cdot (20-x)$. When $x = 0$, $P$ has $y$ coordinate equal to $\frac{11\cdot20}{51} \frac{AP + CP}{AP} = 1 + \frac{CP}{AP}$ = $\frac{51}{20} = 1 + \frac{CP}{AP},$ which equals ${\frac{31}{20}},$ giving an answer of $\boxed{51}.$

 





\textbf{Solution 9 (Menelaus + Ceva's + Angle Bisector Theorem)}

We start by using Menelaus' theorem on $\triangle ABD$ and $EC$.So, we see that $\frac{BC}{DC}\cdot\frac{DM}{AM}\cdot\frac{AE}{EB}=1$.By Angle Bisector theorem, $\frac{BC}{DC}=\frac{31}{11}$, and therefore after plugging in our values we get $\frac{AE}{EB}=\frac{11}{31}$.Then, by Ceva's on the whole figure, we have $\frac{CP}{PA}\cdot\frac{AE}{EB}\cdot\frac{BD}{DC}=1$.Plugging in our values, we get $\frac{CP}{PA}=\frac{31}{20}$, giving an answer of $\boxed{51}$.~ESAOPS

 





\textbf{Solution 10 (Menelaus + ABT)}

By Menelaus' theorem w.r.t $\triangle ACD$ and transversal $\overline{BP}$, we have:

 $\frac{PA}{CP}\cdot\frac{BC}{BD}\cdot\frac{DM}{AM}=1$ $\rightarrow$ $\frac{PA}{CP}\cdot\frac{BC}{BD} = 1$

 Now, we chase the ratio $\frac{BC}{BD}$.

 By the Angle bisector theorem w.r.t $\overline{AD}$ and $\triangle ABC$, we have:

 $\frac{AC}{CD} = \frac{AB}{BD}$ $\rightarrow$ $\frac{CD}{BD} = \frac{11}{20}$ $\rightarrow$ $\frac{BC}{BD} = \frac{31}{20}$

 So, from our Menelaus' equation: $\frac{PA}{CP} = \frac{20}{31}$ $\rightarrow$ $m + n = 20 + 31 = \boxed{051}$.

 


 




	\item (\textit{AIME II 2011}) The sum of the first $2011$ terms of a geometric sequence is $200$. The sum of the first $4022$ terms is $380$. Find the sum of the first $6033$ terms. 


 
 \textbf{Answer: }542



\textbf{Solution}

Since the sum of the first $2011$ terms is $200$, and the sum of the first $4022$ terms is $380$, the sum of the second $2011$ terms is $180$.This is decreasing from the first 2011, so the common ratio is less than one.

 Because it is a geometric sequence and the sum of the first 2011 terms is $200$, second $2011$ is $180$, the ratio of the second $2011$ terms to the first $2011$ terms is $\frac{9}{10}$. Following the same pattern, the sum of the third $2011$ terms is $\frac{9}{10}*180 = 162$.

 Thus, $200+180+162=542$, so the sum of the first $6033$ terms is $\boxed{542}$.

 





\textbf{Solution 2}

Solution by e\_power\_pi\_times\_i

 The sum of the first $2011$ terms can be written as $\dfrac{a_1(1-k^{2011})}{1-k}$, and the first $4022$ terms can be written as $\dfrac{a_1(1-k^{4022})}{1-k}$. Dividing these equations, we get $\dfrac{1-k^{2011}}{1-k^{4022}} = \dfrac{10}{19}$. Noticing that $k^{4022}$ is just the square of $k^{2011}$, we substitute $x = k^{2011}$, so $\dfrac{1}{x+1} = \dfrac{10}{19}$. That means that $k^{2011} = \dfrac{9}{10}$. Since the sum of the first $6033$ terms can be written as $\dfrac{a_1(1-k^{6033})}{1-k}$, dividing gives $\dfrac{1-k^{2011}}{1-k^{6033}}$. Since $k^{6033} = \dfrac{729}{1000}$, plugging all the values in gives $\boxed{542}$.

 





\textbf{Solution 3}

The sum of the first 2011 terms of the sequence is expressible as $a_1 + a_1r + a_1r^2 + a_1r^3$ .... until $a_1r^{2010}$. The sum of the 2011 terms following the first 2011 is expressible as $a_1r^{2011} + a_1r^{2012} + a_1r^{2013}$ .... until $a_1r^{4021}$. Notice that the latter sum of terms can be expressed as $(r^{2011})(a_1 + a_1r + a_1r^2 + a_1r^3...a_1r^{2010})$. We also know that the latter sum of terms can be obtained by subtracting 200 from 380, which then means that $r^{2011} = 9/10$. The terms from 4023 to 6033 can be expressed as $(r^{4022})(a_1 + a_1r + a_1r^2 + a_1r^3...a_1r^{2010})$, which is equivalent to $((9/10)^2)(200) = 162$. Adding 380 and 162 gives the answer of $\boxed{542}$.

 


 





\textbf{Solution 4}

Let $S_n$ be equal to the sum of the first $n$ terms of the geometric sequence. $S_{2011} = 200$ and $S_{4022} = 380$. Let $a$ be the first term and $r$ be the common difference. So $a(\frac{1-r^{2011}}{1-r}) = 200$ and $a(\frac{1-r^{4022}}{1-r}) = 380$. We take the positive difference between the two equations. $a(\frac{r^{2011}-r^{4022}}{r-1}) = 180$. Now, we'll factor out $r^{2011}$ so the equation becomes $ar^{2011}(\frac{1-r^{2011}}{r-1}) = 180$. Divide this equation by the first equation and we get $r^{2011} = 9/10$. We now just need to find the ratio of $S_{6033}$ to $S_{2011}$ multiplied by $S_{2011}$ (It's easy to find the ratio because of common terms). 
 \[\frac{S_{6033}}{S_{2011}} = \frac{\frac{a(1-r^{6033})}{1-r}}{\frac{a(1-r^{2011})}{1-r}} = \frac{1-r^{6033}}{1-r^{2011}} = \frac{1-(\frac{9}{10})^{3}}{1-\frac{9}{10}} = \frac{\frac{271}{1000}}{\frac{1}{10}} = \frac{271}{100}\]Now it's simple, we just need to multiply this value by $S_{2011}$ (which is $200$) and we get our final answer of $\boxed{542}$.

 




	\item (\textit{AIME II 2011}) Define an ordered quadruple of integers $(a, b, c, d)$ as \textit{interesting} if $1 \le a<b<c<d \le 10$, and $a+d>b+c$. How many interesting ordered quadruples are there?


 
 \textbf{Answer: }080



\textbf{Solution 1}

Rearranging the \href{https://artofproblemsolving.com/wiki/index.php?title=Inequality}{inequality} we get $d-c > b-a$. Let $e = 11$, then $(a, b-a, c-b, d-c, e-d)$ is a partition of 11 into 5 positive integers or equivalently:$(a-1, b-a-1, c-b-1, d-c-1, e-d-1)$ is a \href{https://artofproblemsolving.com/wiki/index.php?title=Partition}{partition} of 6 into 5 non-negative integer parts. Via a standard stars and bars argument, the number of ways to partition 6 into 5 non-negative parts is $\binom{6+4}4 = \binom{10}4 = 210$. The interesting quadruples correspond to partitions where the second number is less than the fourth. By symmetry, there are as many partitions where the fourth is less than the second. So, if $N$ is the number of partitions where the second element is equal to the fourth, our answer is $(210-N)/2$.

 We find $N$ as a sum of 4 cases:

 
                \begin{itemize}
                    \item two parts equal to zero, $\binom82 = 28$ ways,\item two parts equal to one, $\binom62 = 15$ ways,\item two parts equal to two, $\binom42 = 6$ ways,\item two parts equal to three, $\binom22 = 1$ way.
                \end{itemize}


                Therefore, $N = 28 + 15 + 6 + 1 = 50$ and our answer is $(210 - 50)/2 = \fbox{080}$

 





\textbf{Solution 2}

Let us consider our quadruple $(a,b,c,d)$ as the following image $xaxbcxxdxx$. The location of the letter $a$, $b$, $c$, $d$ represents its value and $x$ is a place holder. Clearly the quadruple is interesting if there are more place holders between $c$ and $d$ than there are between $a$ and $b$. 

 $0$ holders between $a$ and $b$ means we consider $a$ and $b$ as one unit $ab$ and $c$ as $cx$ yielding $\binom83 = 56$ ways; 

 $1$ holder between $a$ and $b$ means we consider $a$ and $b$ as one unit $axb$ and $c$ as $cxx$ yielding $\binom 63 = 20$ ways; 

 $2$ holders between $a$ and $b$ means we consider $a$ and $b$ as one unit $axxb$ and $c$ as $cxxx$ yielding $\binom43 = 4$ ways.

 Since there cannot be $3$ holders between $a$ and $b$ so our total is $56+20+4=\fbox{080}$.

 





\textbf{Solution 3 (Slightly bashy)}

We first start out when the value of $a=1$. 

 Doing casework, we discover that $d=5,6,7,8,9,10$. We quickly find a pattern.

 Now, doing this for the rest of the values of $a$ and $d$, we see that the answer is simply:

 
$(1)+(2)+(1+3)+(2+4)+(1+3+5)+(2+4+6)+(1)+(2)+(1+3)+(2+4)$$+(1+3+5)+(1)+(2)+(1+3)+(2+4)+(1)+(2)+(1+3)+(1)+(2)+(1)=\boxed{080}$

 





\textbf{Solution 4 (quick)}

Notice that if $a+d>b+c$, then $(11-a)+(11-d)<(11-b)+(11-c)$, so there is a \href{https://artofproblemsolving.com/wiki/index.php?title=Bijection}{bijection} between the number of ordered quadruples with $a+d>b+c$ and the number of ordered quadruples with $a+d<b+c$. 

 Quick counting gives that the number of ordered quadruples with $a+d=b+c$ is 50. To count this, consider our numbers $1, 2, 3, 4, 5, 6, 7, 8, 9, 10$. Notice that if, for example, $a+d=b+c=8$, that the average of $a,d$ and $b,c$ must both be $4$. In this way, there is a symmetry for this case, centered at $4$. If instead, say, $a+d=b+c=7$, an odd number, then there is symmetry with $(a,d);(b,c)$ about $3.5$. Further, the number of cases for each of these centers of symmetry correspond to a triangular number. Eg centered at $2.5$ and $3$ there is $1$ case for each, followed by $3$ cases for $3.5$ and $4$, and so on, until it reaches $5.5$, where there are $10$ possible cases. Then from $6$ to $8.5$, the number of cases for each goes back down all the way to $1$. Adding these all, we have $2(1+1+3+3+6+6)+10=50$.

 Thus the answer is $\frac{\binom{10}{4}-50}{2} = \boxed{080}.$

 





\textbf{Solution 5}

Think about a,b,c,and d as distinct objects, that we must place in 4 of 10 spaces. However, in only 1 of 24 of these combinations, will the placement of these objects satisfy the condition in the problem. So we know the total number of ordered quadruples is $(10*9*8*7/24)=210$ 

 Next, intuitively, the number of quadruples where $a+d>b+c$ is equal to the number of quadruples where $a+d<b+c$. So we need to find the number of quadruples where the two quantities are equal. To do this, all we have to do is consider the cases when $a-d$ ranges from 3 to 9. It would seem natural that a range of 3 would produce 1 option, and a range of 4 would produce 2 options. However, since b and c cannot be equal, a range of 3 or 4 produces 1 option each, a range of 5 or 6 produces 2 options each, a range of 7 or 8 produces 3 options each, and a range of 9 will produce 4 options. In addition, a range of n has 10-n options for combinations of a and d. Multiplying the number of combinations of a and d by the corresponding number of options for b and c gives us 50 total quadruplets where $a+d=b+c$. 

 So the answer will be $\frac{210-50}{2} = \boxed{080}.$

 


 





\textbf{Solution 6}

Let $b = a+x$ and $c = a+x+y$ and $d = a+x+y+z$ for positive integers $x,y,z.$ In order to satisfy the other condition we need $z > x$ so we let $z = x+k.$ Now the only other condition we need to satisfy so $a+2x+y+k \le 10.$ This condition can be transformed into $a+2x+y+k+b  = 11$ for positive $a,x,y,k,b.$ Now we use generating functions to finish. We find the generating function of the whole expression is $(x + x^2 + \cdots)^4 \cdot (x^2+x^4 + \cdots)$ and we are looking for the $x^{11}$ coefficient. This simplifies to finding the $x^5$ coefficient of $(1+x+\cdots)^4 \cdot (1+x^2+\cdots) =\frac{1}{(1-x)^4} \cdot\frac{1}{1+x}.$ Now this expression simplifies to 
 \[\left(\binom{4}{4}+\binom{5}{4} + \cdots +\binom{4+k}{4}x^k\right)(1-x+x^2-x^3 \cdots).\] The $x^5$ coefficient ends up to be $\binom{9}{4} -\binom{8}{4} +\binom{7}{4} -\binom{6}{4} +\binom{5}{4} -\binom{4}{4} = 126 - 70 + 35 - 15 + 5 - 1 = \boxed{080}.$

 





\textbf{Solution 7 (Slightly Slower)}

First, let $a=1$ and $d=10$. If $b=2$, then $c$ can be from $3$ to $8$. If $b=3$, then $4$ to $7$. If $b=4$, then $c$ is between $5$ and $6$. We find a pattern that whenever $b$ increases by $1$, when $a$ and $d$ are stationary, then the possible values of $c$ decrease by 2, unless it gets to zero or negative, in which case that case ends. Counting up, we have $6+4+2=12$ different possibilities when $a=1$ and $d=10$. For $a=1$ and $d=9$, $b=2$, then $c$ can be from $3$ to $7$. If $b=3$, then $c$ can be from $4$ to $6$, and so on. Notice that the possible values for each case of $b$ gets one less than if $d$ were one greater, unless that number is zero, in which it stays zero. We then use this pattern to find all the values: $12+9+6+4+2+1+9+6+4+2+1+6+4+2+1+4+2+1+2+1+1 \Rightarrow  \\ 12\cdot1+9\cdot2+6\cdot3+4\cdot4+2\cdot5+1\cdot6= \\ 12+18+18+16+10+6=\boxed{\textbf{080}}.$

 \textbf{Note} Note that the first value of each of the rows, if we arrange them based on values of $a$, is deleted after each row.

 





\textbf{Solution 8}

Rearranging the equation obtains $b-a<d-c$. Let $a-0=e_0$, $b-a=e_3+e_1$, $c-b=e_2$, $d-c=e_3$, $11-d=e_4$. Add up all of these newly defined equations to obtain $e_0+e_1+e_2+2e_3+e_4=11$. Note that since all $e_n$ were defined to be $e_n\ge1$, to form our stars and bars argument we can let $d_n+1=e_n$ for all $n$. Then we obtain $d_0+d_1+d_2+2d_3+d_4=5$ where $d_n$ is nonnegative. Now, we can move the $2d_3$ term to the other side and perform casework.

 If $2d_3=0$: 5 objects for 4 variables -> $\binom{8}{3}$

 If $2d_3=2$: 3 objects for 4 variables -> $\binom{6}{3}$

 If $2d_3=4$: 1 object for 4 variables -> $\binom{4}{3}$

 Adding all of these cases up, we get $56+20+4=\boxed{080}$ as our requested answer. 

 





\textbf{Solution 9}

Either, $a+d > b+c$, $a+d < b+c$, or $a+d = b+c$. The first two cases are equivalent. The last case is what we need to subtract. 

 Therefore, the number of ways is just $\frac{\binom{10}{4} - N}{2}$ where $N$ is the number of equivalent cases. We have $a+d = b+c$ so $d-c = b-a$. The spacing between the numbers is equivalent. So, denote them as $1 \le a < a+k < c < c+k \le 10$. That implies $1 \le a < c \le 10-k$ where $k$ is an integer. For $k=1$, we have $a+2 \le c \le 9$ so we have 3-9, 4-9, so on until just 9, which is $7+6+dots+1 = 28$ cases. For $k=2$, similar logic yields 15 cases, $k=3$, 6 cases, and $k=4$, 1 case. That is a total 50 cases that are invalidated, so the answer is $\frac{\binom{10}{4} - 50}{2} = \boxed{80}$

 




	\item (\textit{AIME II 2011}) Ed has five identical green marbles, and a large supply of identical red marbles. He arranges the green marbles and some of the red ones in a row and finds that the number of marbles whose right hand neighbor is the same color as themselves is equal to the number of marbles whose right hand neighbor is the other color. An example of such an arrangement is GGRRRGGRG. Let $m$ be the maximum number of red marbles for which such an arrangement is possible, and let $N$ be the number of ways he can arrange the $m+5$ marbles to satisfy the requirement. Find the remainder when $N$ is divided by $1000$.


 
 \textbf{Answer: }003



\textbf{Solution 1}

We are limited by the number of marbles whose right hand neighbor is not the same color as the marble.  By surrounding every green marble with red marbles - RGRGRGRGRGR.  That's 10 "not the same colors" and 0 "same colors."  Now, for every red marble we add, we will add one "same color" pair and keep all 10 "not the same color" pairs.  It follows that we can add 10 more red marbles for a total of $m = 16$.  We can place those ten marbles  in any of 6 "boxes": To the left of the first green marble, to the right of the first but left of the second, etc. up until to the right of the last.  This is a stars-and-bars problem, the solution of which can be found as $\binom{n+k}{k}$ where n is the number of stars and k is the number of bars.  There are 10 stars (The unassigned Rs, since each "box" must contain at least one, are not counted here) and 5 "bars," the green marbles.  So the answer is $\binom{15}{5} = 3003$, take the remainder when divided by 1000 to get the answer: $\boxed{003}$.

 
Rewording:

 There are three states for every marble: the one to the right is different, the one to the right is the same, and there is no marble to its right. Since there are unlimited red marbles, the second value can be created infinitely and is therefore worthless. However, the first value can only be created with a combination green marbles, thus finite, and worthy of our interest. As such, to maximize the amount of red marbles, we need to maximize the appearances of the first value. The "minimal" of this strategy is to simply separate every green marble with a red - RGRGRGRGRGR. There are ten appearances of the first value, so we must add a number of red marbles to add appearances of the second value. With quick testing we can conclude the need of nine new red marbles that can be placed in six defined places that are divided by the green marbles. This is followed by stars and bars which gives us $\binom{15}{5} = 3003$, $\boxed{003}$.

 
-jackshi2006

 





\textbf{Solution 2}

Begin with the above solution's reasoning to find that the ideal sequence is RGRGRGRGRGR + 10 Rs. To count the number of arrangements where the Gs are separate, group a G and an R together (GR) which will be rearranged/counted as one. However this also introduces the case of having GRGR.... which results in too few number of marbles whose right-hand neighbor is different, so we fix an R in the beginning of our sequence which also removes it from our counting calculation. Finally we are looking at the arrangement of 5 (GR)s + 10 Rs, all indistinguishable, which is $\frac{15!}{5!10!}=3003\Longrightarrow\boxed{003}$.

 





\textbf{Solution 3}

Let a collection of consecutive letters be a "blob."  Let there be n blobs.

 The marbles such that the marble to their right has a different color are the ones that are the last marble of a "blob."  The only exception to this is the very last marble in this list. There are \( n-1 \) of these.  

 The marbles such that the marble to their right has the same color are any marbles which are not at the end of a blob.  The only exception to this is the last marble, which does not go in any category.  

 The problem states that the number of marbles like this is equal to the number of marbles with a different color.  Therefore, there are \( (n-1) + (n-1) + 1 \) marbles in total (the last \( +1 \) is for the last marble at the end).  We know there are \( m+5 \) marbles, so \( m+5 = 2n-1 \).  

 Solving for \( n \), we get \( n = \frac{m}{2} + 3 \).  

 Notice that \( n \leq 11 \) because the "best-case scenario" is that each \( G \) forms its own blob, and there is a red blob in between any two \( G \)'s and a red blob on both ends.  Here is how it should look:  

 <red blob> G <red blob> G <red blob> G <red blob> G <red blob> G <red blob>.

 
In this case, there are \( 11 \) blobs.  

 Thus, \( \frac{m}{2} + 3 \leq 11 \implies m \leq 16 \).  

 This is now a stars and bars problem. For \( m = 16 \), we need to create the optimal solution: \( 11 \) blobs.  Since every <red blob> must have at least one red marble, we can just give each blob one red at the start, leaving us with \( 16-6 = 10 \) red marbles left to separate into \( 6 \) groups.  By stars and bars, this is \( \binom{15}{5} = 3003 \).  

 Thus, the answer is \( \boxed{003} \).

 





\textbf{Solution 4 (VERY SIMILAR TO 3)}

Let a collection of consecutive letters be a "blob."  Let the blobs have sizes s\_1, s\_2, ... s\_n given the existence of n blobs.

 The marbles such that the marble to their right has a different color are the ones that are the last marble of a "blob."  The only exception to this is the very last marble in this list. There are \( n-1 \) of these.  

 The marbles such that the marble to their right has the same color is every marble in a blob but the last, which is 
 \[\sum_{i=1}^{n} (s_i - 1) = \left(\sum_{i=1}^{n} s_i\right) - n = m + 5 - n.\]

 


 Equating our two expression $n-1 = m+5-n \rightarrow 2n=m+6.$

 
Solving for \( n \), we get \( n = \frac{m}{2} + 3 \).  

 Notice that \( n \leq 11 \) because the "best-case scenario" is that each \( G \) forms its own blob, and there is a red blob in between any two \( G \)'s and a red blob on both ends.  Here is how it should look:  

 <red blob> G <red blob> G <red blob> G <red blob> G <red blob> G <red blob>.

 
In this case, there are \( 11 \) blobs.  

 Thus, \( \frac{m}{2} + 3 \leq 11 \implies m \leq 16 \).  

 This is now a stars and bars problem. For \( m = 16 \), we need to create the optimal solution: \( 11 \) blobs.  Since every <red blob> must have at least one red marble, we can just give each blob one red at the start, leaving us with \( 16-6 = 10 \) red marbles left to separate into \( 6 \) groups.  By stars and bars, this is \( \binom{15}{5} = 3003 \).  

 Thus, the answer is \( \boxed{003} \).

 




	\item (\textit{AIME II 2011}) Let $z_1,z_2,z_3,\dots,z_{12}$ be the 12 zeroes of the polynomial $z^{12}-2^{36}$. For each $j$, let $w_j$ be one of $z_j$ or $i z_j$. Then the maximum possible value of the real part of $\sum_{j=1}^{12} w_j$ can be written as $m+\sqrt{n}$ where $m$ and $n$ are positive integers. Find $m+n$.


 
 \textbf{Answer: }784



\textbf{Solution}


                    \begin{figure}[H]
                        \centering
                        \includegraphics[width=0.5\textwidth]{Solution/AIME_II_2011_8_Solution_Figure_1.png}
                    \end{figure}
                    The twelve dots above represent the $12$ roots of the equation $z^{12}-2^{36}=0$. If we write $z=a+bi$, then the real part of $z$ is $a$ and the real part of $iz$ is $-b$. The blue dots represent those roots $z$ for which the real part of $z$ is greater than the real part of $iz$, and the red dots represent those roots $z$ for which the real part of $iz$ is greater than the real part of $z$. Now, the sum of the real parts of the blue dots is easily seen to be $8+16\cos\frac{\pi}{6}=8+8\sqrt{3}$ and the negative of the sum of the imaginary parts of the red dots is easily seen to also be $8+8\sqrt{3}$. Hence our desired sum is $16+16\sqrt{3}=16+\sqrt{768}$, giving the answer $\boxed{784}.$

 





\textbf{Solution }

As a small note, we could factor the equation as $z^{12} = 2^{36} \implies (\frac{z}{8})^{12} = 1$ in which these are just the 12th roots of unity for $\frac{z}{8}$ which might be easier to work with. Then we just add up the optimal real parts as in solution 1 and then multiply by 8 at the end.

 





\textbf{Solution }

Note that $\sin(x) = \sin(x + \pi/2 - \pi/2) = -\cos(x + \pi/2)$. When summing the real part of $8i(\cos(x) + i\sin(x))$ is $-8\sin(x) = 8\cos(x + \pi/2)$.

 So, translate all the red dots by $90^{\circ}$ and sum the cosines of angles of the blue and (rotated) red dots.

 
                    \begin{figure}[H]
                        \centering
                        \includegraphics[width=0.5\textwidth]{Solution/AIME_II_2011_8_Solution_Figure_2.png}
                    \end{figure}
                    

 Then, it's easy to see that points 2, 3, and 4 cancel out with 5, 6, and 7, so you're left with $8 \cdot (\cos \pi/6 + \cos -\pi/6 + \cos 0)$ times two~\href{https://artofproblemsolving.com/wiki/index.php?title=User:Crazyvideogamez}{CrazyVideoGamez}

 





\textbf{Solution 2}

The equation $z^{12}-2^{36}=0$ can be factored as follows:
 \[(z^6-2^{18})(z^6+2^{18})=0\]
 \[(z^2-2^6)(z^2+2^6)({(z^2+2^6)}^2-z^2\cdot2^6)({(z^2-2^6)}^2+z^2\cdot2^6)=0\]
 \[(z^2-2^6)(z^2+2^6)(z^2+2^6-z\cdot2^3)(z^2+2^6+z\cdot2^3) (z^2-2^6-iz\cdot2^3)(z^2-2^6+i z\cdot2^3)=0\]

 Since this is a 12th degree equation, there are 12 roots. Also, since each term in the equation is even, the positive or negative value of each root is another root. That would mean there are 6 roots that can be multiplied by $-1$ and since we have 6 factors, that’s 1 root per factor. We just need to solve for $z$ in each factor and pick whether or not to multiply by $i$ and $-1$ for each one depending on the one that yields the highest real value. After that process, we get $8+8+2((4\sqrt{3}+4)+(4\sqrt{3}-4))$ Adding the values up yields $16+16\sqrt{3}$, or $16+\sqrt{768}$, and $16+768=\boxed{784}$.

 -Solution by Someonenumber011.

 





\textbf{Solution 3}

Clearly, the roots are: $2^3*(\cos{\frac{k\pi}{12}}+i\sin{\frac{k\pi}{12}}), k\in{0, 1, 2, 3, 4, 5, 6, 7, 8, 9, 10, 11}$. 

 Now, realize for $z=a+bi$, $\operatorname{Re}(iz)=-b$. $\operatorname{Re}(z)=a$. $\operatorname{Re}(z)<\operatorname{Re}(iz)$ is true when $a<-b$.

 This means:

 When $a>0$, $b<-a<0$. 

 When $a<0$, $0<b<-a$.

 For the 12 roots of the polynomial in the original equation, $8\cos{k\pi/12}=\operatorname{Re}(z), k\in{0, 1, 2, 3, 4, 5, 6, 7, 8, 9, 10, 11}$. $-8\sin{k\pi/12}=\operatorname{Im}(iz), k\in{0, 1, 2, 3, 4, 5, 6, 7, 8, 9, 10, 11}$.

 So $8\cos{k\pi/12}<-8\sin{k\pi/12}$. $-\cos{k\pi/12}>\sin{k\pi/12}$. This can be easily true for roots that are in the 3rd quadrant in the complex plane. 

 This cannot be true for roots in the 1st quadrant because that would yield a negative number bigger than a positive one.

 Consider the roots in the 2nd and 4th quadrants. Calculate the roots, choose, and then add the ones up. You will get $\boxed{784}$.

 





\textbf{Solution 4 (Brute Force)}

Use De Moivre's Theorem to brute force all the roots out. Then choose the greater value of $\operatorname{Re}(z), \operatorname{Re}(iz)$. After adding everything up, you get $\boxed{784}$.

 




	\item (\textit{AIME II 2011}) Let $x_1$, $x_2$, $\dots$, $x_6$ be nonnegative real numbers such that $x_1 + x_2 + x_3 + x_4 + x_5 + x_6 = 1$, and $x_1x_3x_5 + x_2x_4x_6 \ge {\frac{1}{540}}$. Let $p$ and $q$ be relatively prime positive integers such that $\frac{p}{q}$ is the maximum possible value of $x_1x_2x_3 + x_2x_3x_4 + x_3x_4x_5 + x_4x_5x_6 + x_5x_6x_1 + x_6x_1x_2$. Find $p + q$.


 
 \textbf{Answer: }559



\textbf{Solution 1}

Note that neither the constraint nor the expression we need to maximize involves products $x_i x_j$ with $i \equiv j \pmod 3$. Factoring out say $x_1$ and $x_4$ we see that the constraint is $x_1(x_3x_5) + x_4(x_2x_6) \ge {\frac1{540}}$, while the expression we want to maximize is $x_1(x_2x_3 + x_5x_6 + x_6x_2) + x_4(x_2x_3 + x_5x_6 + x_3x_5)$. Adding the left side of the constraint to the expression, we get: $(x_1 + x_4)(x_2x_3 + x_5x_6 + x_6x_2 + x_3x_5) = (x_1 + x_4)(x_2 + x_5)(x_3 + x_6)$. This new expression is the product of three non-negative terms whose sum is equal to 1. By AM-GM this product is at most $\frac1{27}$. Since we have added at least $\frac{1}{540}$ the desired maximum is at most $\frac1{27} - \frac1{540} =\frac{19}{540}$. It is easy to see that this upper bound can in fact be achieved by ensuring that the constraint expression is equal to $\frac1{540}$ with $x_1 + x_4 = x_2 + x_5 = x_3 + x_6 =\frac13$—for example, by choosing $x_1$ and $x_2$ small enough—so our answer is $540 + 19 = \fbox{559}.$

 An example is:
 \begin{align*} x_3 &= x_6 = \frac16 \\ x_1 &= x_2 = \frac{5 - \sqrt{20}}{30} \\ x_5 &= x_4 = \frac{5 + \sqrt{20}}{30} \end{align*}


 Another example is 
 \begin{align*} x_1 = x_3 = \frac{1}{3} \\ x_2 = \frac{19}{60}, \ x_5 = \frac{1}{60} \\ x_4 &= x_6 = 0 \end{align*}


 





\textbf{Solution 2 (Not legit)}

There's a symmetry between $x_1, x_3, x_5$ and $x_2,x_4,x_6$. Therefore, a good guess is that $a = x_1 = x_3 = x_5$ and $b = x_2 = x_4 = x_6$, at which point we know that $a+b = 1/3$, $a^3+b^3 \geq 1/540$, and we are trying to maximize $3a^2b+3ab^2$. Then,

 
 \[3a^2b+3ab^2 = (a+b)^3-a^3-b^3 \leq \frac{1}{27} - \frac{1}{540} = \boxed{\frac{19}{540}}\] which is the answer.

 This solution is extremely lucky; if you attempt to solve for $a$ and $b$ you receive complex answers (which contradict the problem statement), but the final answer is correct.

 




	\item (\textit{AIME II 2011}) A circle with center $O$ has radius 25. Chord $\overline{AB}$ of length 30 and chord $\overline{CD}$ of length 14 intersect at point $P$. The distance between the midpoints of the two chords is 12. The quantity $OP^2$ can be represented as $\frac{m}{n}$, where $m$ and $n$ are relatively prime positive integers. Find the remainder when $m + n$ is divided by 1000.


 
 \textbf{Answer: }057



\textbf{Solution 1}

Let $E$ and $F$ be the midpoints of $\overline{AB}$ and $\overline{CD}$, respectively, such that $\overline{BE}$ intersects $\overline{CF}$.

 Since $E$ and $F$ are midpoints, $BE = 15$ and $CF = 7$.

 $B$ and $C$ are located on the circumference of the circle, so $OB = OC = 25$.

 The line through the midpoint of a chord of a circle and the center of that circle is perpendicular to that chord, so $\triangle OEB$ and $\triangle OFC$ are right triangles (with $\angle OEB$ and $\angle OFC$ being the right angles).  By the Pythagorean Theorem, $OE = \sqrt{25^2 - 15^2} = 20$, and $OF = \sqrt{25^2 - 7^2} = 24$.

 Let $x$, $a$, and $b$ be lengths $OP$, $EP$, and $FP$, respectively.  OEP and OFP are also right triangles, so $x^2 = a^2 + 20^2 \to a^2 = x^2 - 400$, and $x^2 = b^2 + 24^2 \to b^2 = x^2 - 576$

 We are given that $EF$ has length 12, so, using the \href{https://artofproblemsolving.com/wiki/index.php?title=Law\_of\_Cosines}{Law of Cosines} with $\triangle EPF$:

 $12^2 = a^2 + b^2 - 2ab \cos (\angle EPF) = a^2 + b^2 - 2ab \cos (\angle EPO + \angle FPO)$

 Substituting for $a$ and $b$, and applying the Cosine of Sum formula:

 $144 = (x^2 - 400) + (x^2 - 576) - 2 \sqrt{x^2 - 400} \sqrt{x^2 - 576} \left( \cos \angle EPO \cos \angle FPO - \sin \angle EPO \sin \angle FPO \right)$

 $\angle EPO$ and $\angle FPO$ are acute angles in right triangles, so substitute opposite/hypotenuse for sines and adjacent/hypotenuse for cosines:

 $144 = 2x^2 - 976 - 2 \sqrt{(x^2 - 400)(x^2 - 576)} \left(\frac{\sqrt{x^2 - 400}}{x} \frac{\sqrt{x^2 - 576}}{x} - \frac{20}{x} \frac{24}{x} \right)$

 Combine terms and multiply both sides by $x^2$: $144 x^2 = 2 x^4 - 976 x^2 - 2 (x^2 - 400) (x^2 - 576) + 960  \sqrt{(x^2 - 400)(x^2 - 576)}$

 Combine terms again, and divide both sides by 64: $13 x^2 = 7200 - 15 \sqrt{x^4 - 976 x^2 + 230400}$

 Square both sides: $169 x^4 - 187000 x^2 + 51,840,000 = 225 x^4 - 219600 x^2 + 51,840,000$

 This reduces to $x^2 = \frac{4050}{7} = (OP)^2$; $4050 + 7 \equiv \boxed{057} \pmod{1000}$.

 





\textbf{Solution 2 - Fastest}

We begin as in the first solution.  Once we see that $\triangle EOF$ has side lengths $12$, $20$, and $24$, we can compute its area with Heron's formula:

 
 \[K = \sqrt{s(s-a)(s-b)(s-c)} = \sqrt{28\cdot 16\cdot 8\cdot 4} = 32\sqrt{14}.\]

 Thus, the circumradius of triangle $\triangle EOF$ is $R = \frac{abc}{4K} = \frac{45}{\sqrt{14}}$.  Looking at $EPFO$, we see that $\angle OEP = \angle OFP = 90^\circ$, which makes it a cyclic quadrilateral. This means $\triangle EOF$'s circumcircle and $EPFO$'s inscribed circle are the same.

 Since $EPFO$ is \href{https://artofproblemsolving.com/wiki/index.php?title=Cyclic}{cyclic} with diameter $OP$, we have $OP = 2R = \frac{90}{\sqrt{14}}$, so $OP^2 = \frac{4050}{7}$ and the answer is $\boxed{057}$.

 





\textbf{Solution 3}

We begin as the first solution have $OE=20$ and $OF=24$. Because $\angle PEO+\angle PFO=180^{\circ}$, Quadrilateral $EPFO$ is inscribed in a Circle. Assume point $I$ is the center of this circle. 

 $\because \angle OEP=90^{\circ}$

 $\therefore$ point $I$ is on $OP$

 Link $EI$ and $FI$, Made line $IK\bot EF$, then $\angle EIK=\angle EOF$

 On the other hand, $\cos\angle EOF=\frac{EO^2+OF^2-EF^2}{2\cdot EO\cdot OF}=\frac{13}{15}=\cos\angle  EIK$

 $\sin\angle EOF=\sin\angle EIK=\sqrt{1-\frac{13^2}{15^2}}=\frac{2\sqrt{14}}{15}$

 As a result, $IE=IO=\frac{45}{\sqrt 14}$

 Therefore, $OP^2=4\cdot \frac{45^2}{14}=\frac{4050}{7}.$

 As a result, $m+n=4057\equiv \boxed{057}\pmod{1000}$

 





\textbf{Solution 4}

Let $OP=x$.

 Proceed as the first solution in finding that quadrilateral $EPFO$ has side lengths $OE=20$, $OF=24$, $EP=\sqrt{x^2-20^2}$, and $PF=\sqrt{x^2-24^2}$, and diagonals $OP=x$ and $EF=12$.

 We note that quadrilateral $EPFO$ is cyclic and use Ptolemy's theorem to solve for $x$:

 
 \[20\cdot \sqrt{x^2-24^2} + 12\cdot x = 24\cdot \sqrt{x^2-20^2}\]

 Solving, we have $x^2=\frac{4050}{7}$ so the answer is $\boxed{057}$.

 -Solution by ilovebigmeat

 





\textbf{Solution 5 (Quick Angle Solution)}

Let $M$ be the midpoint of $AB$ and $N$ of $CD$. As $\angle OMP = \angle ONP$, quadrilateral $OMPN$ is cyclic with diameter $OP$. By Cyclic quadrilaterals note that $\angle MPO = \angle MNO$.

 The area of $\triangle MNP$ can be computed by Herons as 
 \[[MNO] = \sqrt{s(s-a)(s-b)(s-c)} = \sqrt{28\cdot 16\cdot 8\cdot 4} = 32\sqrt{14}.\] The area is also $\frac{1}{2}ON \cdot MN \sin{\angle MNO}$. Therefore, 
 \begin{align*} \sin{\angle MNO} &= \frac{2[MNO]}{ON \cdot MN}  \\ &= \frac{2}{9}\sqrt{14} \\ \sin{\angle MNO} &= \frac{OM}{OP} \\ &= \frac{2}{9}\sqrt{14} \\ OP &= \frac{90\sqrt{14}}{14} \\ OP^2 &= \frac{4050}{7} \implies \boxed{057}. \end{align*}
 

 





\textbf{Solution 6}

Define $M$ and $N$ as the midpoints of $AB$ and $CD$, respectively. Because $\angle OMP = \angle ONP = 90^{\circ}$, we have that $ONPM$ is a cyclic quadrilateral. Hence, $\angle PNM = \angle POM.$ Then, let these two angles be denoted as $\alpha$. Now, assume WLOG that $PD = x < 7$ and $PB = y < 15$ (We can do this because one of $PD$ or $PC$ must be less than 7, and similarly for $PB$ and $PA$). Then, by Power of a Point on P with respect to the circle with center $O$, we have that 
 \[(14-x)x = (30-y)y\]
 \[(7-x)^{2}+176=(15-y)^{2}.\]Then, let $z = (7-x)^{2}$. From Law of Cosines on $\triangle NMP$, we have that 
 \[\textrm{cos } \angle MNP = \frac{NP^{2}+MN^{2}-MP^{2}}{2 \cdot NP \cdot MN}\]
 \[\textrm{cos } \alpha  = \frac{(7-x)^{2} + 12^{2} - (14-x)^{2}}{24 \cdot (7-x)}.\]Plugging in $z$ in gives
 \[\textrm{cos } \alpha  = \frac{-32}{24 \cdot \sqrt{z}}\]
 \[\textrm{cos } \alpha  = \frac{-4}{3\sqrt{z}}\]
 \[\textrm{cos }^{2} \alpha  = \frac{16}{9z}.\]Hence,
 \[\textrm{tan }^{2} \alpha  = \frac{\frac{9z-16}{9z}}{\frac{16}{9z}} = \frac{9z-16}{16}.\]Then, we also know that 
 \[\textrm{tan } \alpha  = \textrm{tan } \angle MOP  = \frac{MP}{OM} = \frac{14-y}{20}.\]Squaring this, we get
 \[\textrm{tan }^{2} \alpha  = \frac{z+176}{400}.\]Equating our expressions for $z$, we get$\frac{z+176}{400} = \frac{9z-16}{16}.$Solving gives us that$z = \frac{18}{7}$.Since $\angle ONP = 90^{\circ}$, from the Pythagorean Theorem, $OP^{2} = ON^{2}+PN^{2} = 25^{2}-7^{2} + z = 576+z = \frac{4050}{7}$,and thus the answer is $4050+7 = 4057$, which when divided by a thousand leaves a remainder of $\boxed{57}.$

 -Mr.Sharkman

 Note: my solution was very long and tedious. It was definitely the least elegant solution. The only thing I like about it is it contains no quadratic equations (unless you count LoC).

 





\textbf{Solution 7 Analytic Geometry}


                    \begin{figure}[H]
                        \centering
                        \includegraphics[width=0.5\textwidth]{Solution/AIME_II_2011_10_Solution_Figure_1.png}
                    \end{figure}
                    

 Let $E$ and $F$ be the midpoints of $\overline{AB}$ and $\overline{CD}$, respectively, such that $\overline{BE}$ intersects $\overline{CF}$.

 Since $E$ and $F$ are midpoints, $BE = 15$ and $CF = 7$.

 $B$ and $C$ are located on the circumference of the circle, so $OB = OC = 25$.

 Since $\overline{OE}\perp \overline{AB}$ and $\overline{OF}\perp \overline{CD}$,$OE = \sqrt{OB^2-BE^2}=20$ and $OF = \sqrt{OC^2-OF^2}=24$

 With law of cosines, $\cos \angle EOF = \frac{OE^2+OF^2-EF^2}{2\cdot OE\cdot OF} = \frac{13}{15}$ 

 Since $EF < OF$, $\angle EOF$ is acute angle. $\sin \angle EOF = \sqrt{1-\cos^2 \angle EOF} = \frac{\sqrt{56}}{15}$ and $\tan \angle EOF = \frac{\sqrt{56}}{13}$

 Let $\overline{OF}$ line be $x$ axis. 

 Line $\overline{DC}$ equation is $x = OF$. 

 Since line $\overline{AB}$ passes point $E$ and perpendicular to $\overline{OD}$, its equation is $y - E_y = -\frac{1}{\tan \angle EOF} (x - E_x)$

 where $E_x = OE\cos{\angle EOF}$ , $E_y = OE\sin{\angle EOF}$

 Since $P$ is the intersection of $\overline{AB}$ and $\overline{CD}$,

 $P_x = OF = 24$

 $P_y = E_y -\frac{1}{\tan \angle EOF} (OF - E_x) = - \frac{3\sqrt{14}}{7}$ (Negative means point $P$ is between point $F$ and $C$)

 $OP^2 = P_x^2 + P_y^2 = \frac{4050}{7}$ and the answer is $\boxed{057}$. 

 Note: if $EF$ was longer, point $P$ would be between point $D$ and $F$. Then, $OP$ would be the diagonal of quadrilateral $OEPF$ not the side. To apply Ptolemy's theorem like solution 4, it is critical to know whether $OP$ is the diagonal or side of quadrilateral. Equation for wrong case cannot be solved. For example,

 
                    \begin{figure}[H]
                        \centering
                        \includegraphics[width=0.5\textwidth]{Solution/AIME_II_2011_10_Solution_Figure_2.png}
                    \end{figure}
                    

 





\textbf{Solution 8 (Analytic- Can use complex numbers or rotation matrix)}


 We work in the complex plane where \( O = 0 \). Let \( M_1, M_2 \), be midpoints of chords \( AB \), and \( CD \) respectively. \( AB=30 \rightarrow AM=15 \). Since \( AM=15 \), and \( OA=25 \), \( OM_1 = 20 \) by the Pythagorean theorem. WLOG, Let \( M_1 = 20 + 0i \).  Think about the locus of all points \( M_2 \). By similar logic as above it is any \( M_2 \), such that \( OM_2=24 \). Note that \( M_1M_2 = 12 \). Therefore by LoC on \( \triangle OM_1M_2 \), if \( \angle O = \theta \), \( \cos\theta = \frac{13}{15} \rightarrow \sin\theta = \frac{2\sqrt{14}}{15} \).  Note that a rotation by \( \theta \) and a scaling of \( \frac{24}{20} = \frac{6}{5} \) transforms \( AB \) to \( DC \).  This transformation = \( \text{cis}\theta \cdot \frac{6}{5} = \frac{2}{5} \cdot \frac{13 + 2\sqrt{14}i}{5} \).  Applying the transformation we on \( A,B \) respectively we find: (We are back in Cartesian coordinates now)  \[ D = \frac{2}{5} \cdot (52 - 6\sqrt{14}, 39 + 8\sqrt{14}), \] \[ C = \frac{2}{5} \cdot (52 + 6\sqrt{14}, -39 + 8\sqrt{14}). \]  To make bashing easier first ignore the \( \frac{2}{5} \).   e.g. let \( D = 0.4D' \) and \( C = 0.4C' \) and first find the equation of the line passing through \( C' \) and \( D' \).   After bashing a bit, we get  \[ y = \frac{-13}{2\sqrt{14}} \cdot x + \frac{225}{7}\sqrt{14}. \]  Now to account for the \( \frac{2}{5} \) thing we say the actual line is this: \[ y = \frac{-13}{2\sqrt{14}} \cdot x + \frac{2}{5} \cdot \frac{225}{7} \sqrt{14}, \] \[ y = \frac{-13}{2\sqrt{14}} \cdot x + \frac{90}{7} \sqrt{14}. \]  \( P \) is the intersection of \( CD \) and \( AB \). We have the equation of \( CD \), and \( AB \) is simply \( x = 20 \), so   Letting \( x = 20 \) we find  \[ y = \frac{25\sqrt{14}}{7}. \]   \[ OP^2 = 20^2 + \left( \frac{25\sqrt{14}}{7} \right)^2 = 5^2 \cdot \left( 4^2 + \left( \frac{5\sqrt{2}}{\sqrt{7}} \right)^2 \right) = 25 \cdot \frac{810}{7} = \frac{4050}{7}. \] Thus, the answer is \( \boxed{057} \).

 





\textbf{Solution 9}

Let the midpoint of $AB$ be $M$, and the midpoint of $CD$ be $N$. Because $ON$ is a perpendicular bisector of $CD$, we know that $CN=7$ and $OC=25$ (the radius of the circle). Therefore, by the Pythagorean Theorem, $ON=24$.

 Now that we have all the side lengths of $\triangle{OMN}$, we use Law of Cosines to find $\angle{OMN}$.\begin{align*}20^2+12^2-(2)(20)(12)(\cos\angle{OMN})&=24^2\\544-480\cos\angle{OMN}&=576\\\cos\angle{OMN}&=-\dfrac{1}{15}.\end{align*}
Because $\angle{OMP}$ is right, we know that $\angle{PMN}=\angle{OMN}-90^\circ$. 

 We then use the cosine difference identity to find that (note that we use $\cos 90^\circ =0$ and $\sin 90^\circ = 1$)\begin{align*}\cos\angle{PMN}&=\cos(\angle{OMN}-90^\circ)\\&=\cos\angle{OMN}\cos 90^\circ + \sin\angle{OMN}\sin 90^\circ\\&=\sin{OMN}.\end{align*}
Using the identity $\cos^2 \theta+\sin^2 \theta=1$, and since it is an acute angle it must be positive, we find that $\cos\angle{PMN}=\dfrac{4\sqrt{14}}{15}$.

 Let $PM=x$, and $PN=y$. Using Law of Sines on $\triangle{PMN}$, we have the equation
 \[x^2+144-(2)(x)(12)(\cos\angle{PMN})=y^2.\]Simplifying, we get
 \[x^2-y^2+144=\dfrac{96x\sqrt{14}}{15}.\]Using Power of a Point on point $P$, we find that $(15-x)(15+x)=(7-y)(7+y)$, or $x^2-y^2=176$. Substituting this in and simplifying, we find that$x=\dfrac{50}{\sqrt{14}}$.Finally, using the Pythagorean Theorem on $\triangle{OMP}$, we find that
 \[OP^2=400+x^2=400+\dfrac{2500}{14}=\dfrac{2800+1250}{7}=\dfrac{4050}{7}.\]Thus, the answer to the problem is $4050+7=4057\rightarrow \boxed{057}.$

 - The\_Other\_Guy

 





\textbf{Solution 10}

Place the circle of radius $25$ at the origin, $O=(0,0)$. Let chord $\overline{AB}$ be vertical with midpoint $M=(-20,0)$, so that $AB$ has length $30$. Then the endpoints are $A=(-20,-15), \quad B=(-20,15)$.

 Let chord $\overline{CD}$ have length $14$, so its half-length is $7$. The distance from the origin to the midpoint of $\overline{CD}$ is $\sqrt{25^2 - 7^2} = \sqrt{625 - 49} = 24.$

 Let the midpoint of $\overline{CD}$ be $N=(x,y)$. Since the distance between the midpoints of the chords is $12$, we have $(x+20)^2 + y^2 = 12^2 = 144$, and from the distance to the origin: $x^2 + y^2 = 24^2 = 576$.

 Subtracting these equations eliminates $y^2$ and gives a simple linear equation:

 $(x+20)^2 + y^2 - (x^2 + y^2) = 144 - 576 \Rightarrow x^2 + 40x + 400 - x^2 = -432 \Rightarrow 40x = -832 \Rightarrow x = -\frac{104}{5}.$

 Then, $y^2 = 144 - (x+20)^2 = 144 - \left(-\frac{4}{5}\right)^2 = \frac{3600}{25} - \frac{16}{25} = \frac{3584}{25} \Rightarrow y = \frac{16\sqrt{14}}{5}.$

 Thus the midpoint of $\overline{CD}$ is $N = \left(-\frac{104}{5}, \frac{16\sqrt{14}}{5}\right)$.

 The slope of $\overline{CD}$ is perpendicular to the radius at its midpoint: $\text{slope} = \frac{-(-104/5)}{16\sqrt{14}/5} = \frac{104/5}{16\sqrt{14}/5} = \frac{13}{2\sqrt{14}}$.

 Equation of the line through $N$:$y - \frac{16\sqrt{14}}{5} = \frac{13}{2\sqrt{14}} \left(x + \frac{104}{5}\right)$.

 Intersecting with $x=-20$ (the vertical line for AB): $y - \frac{16\sqrt{14}}{5} = \frac{13}{2\sqrt{14}} \cdot \frac{4}{5} \Rightarrow y = \frac{16\sqrt{14}}{5} + \frac{26}{5\sqrt{14}} = \frac{50}{\sqrt{14}}$.

 So the intersection point $P$ is $P = \left(-20, \frac{50}{\sqrt{14}}\right)$, and $OP^2 = (-20)^2 + \left(\frac{50}{\sqrt{14}}\right)^2 = 400 + \frac{2500}{14} = \frac{4050}{7}$.

 Thus $m+n = 4050 + 7 = 4057$, and the remainder when divided by $1000$ is $\boxed{57}.$

 




	\item (\textit{AIME II 2011}) Let $M_n$ be the $n \times n$ matrix with entries as follows: for $1 \le i \le n$, $m_{i,i} = 10$; for $1 \le i \le n - 1$, $m_{i+1,i} = m_{i,i+1} = 3$; all other entries in $M_n$ are zero. Let $D_n$ be the determinant of matrix $M_n$. Then $\sum_{n=1}^{\infty} \frac{1}{8D_n+1}$ can be represented as $\frac{p}{q}$, where $p$ and $q$ are relatively prime positive integers. Find $p + q$. 


 Note: The determinant of the $1 \times 1$ matrix $[a]$ is $a$, and the determinant of the $2 \times 2$ matrix $\left[ {\begin{array}{cc}  a & b  \\  c & d  \\  \end{array} } \right] = ad - bc$; for $n \ge 2$, the determinant of an $n \times n$ matrix with first row or first column $a_1$ $a_2$ $a_3$ $\dots$ $a_n$ is equal to $a_1C_1 - a_2C_2 + a_3C_3 - \dots + (-1)^{n+1}a_nC_n$, where $C_i$ is the determinant of the $(n - 1) \times (n - 1)$ matrix formed by eliminating the row and column containing $a_i$.


 
 \textbf{Answer: }073



\textbf{Solution}


 \[D_{1}=\begin{vmatrix}  10    \end{vmatrix} = 10, \quad D_{2}=\begin{vmatrix} 10 & 3  \\  3 & 10  \\ \end{vmatrix} =(10)(10) - (3)(3) = 91, \quad D_{3}=\begin{vmatrix}  10 & 3 & 0  \\  3 & 10 & 3  \\  0 & 3 & 10  \\  \end{vmatrix}.\]Using the expansionary/recursive definition of determinants (also stated in the problem):

 $D_{3}=\left| {\begin{array}{ccc}  10 & 3 & 0  \\  3 & 10 & 3  \\  0 & 3 & 10  \\  \end{array} } \right|=10\left| {\begin{array}{cc}  10 & 3  \\  3 & 10  \\  \end{array} } \right| - 3\left| {\begin{array}{cc}  3 & 3  \\  0 & 10  \\  \end{array} } \right| + 0\left| {\begin{array}{cc}  3 & 10  \\  0 & 3  \\  \end{array} } \right| = 10D_{2} - 9D_{1} = 820$

 This pattern repeats because the first element in the first row of $M_{n}$ is always 10, the second element is always 3, and the rest are always 0. The ten element directly expands to $10D_{n-1}$. The three element expands to 3 times the determinant of the the matrix formed from omitting the second column and first row from the original matrix. Call this matrix $X_{n}$. $X_{n}$ has a first column entirely of zeros except for the first element, which is a three. A property of matrices is that the determinant can be expanded over the rows instead of the columns (still using the recursive definition as given in the problem), and the determinant found will still be the same. Thus, expanding over this first column yields $3D_{n-2} + 0(\text{other things})=3D_{n-2}$. Thus, the $3 \det(X_{n})$ expression turns into $9D_{n-2}$. Thus, the equation $D_{n}=10D_{n-1}-9D_{n-2}$ holds for all n > 2.

 This equation can be rewritten as $D_{n}=10(D_{n-1}-D_{n-2}) + D_{n-2}$. This version of the equation involves the difference of successive terms of a recursive sequence. Calculating $D_{0}$ backwards from the recursive formula and $D_{4}$ from the formula yields $D_{0}=1, D_{4}=7381$. Examining the differences between successive terms, a pattern emerges. $D_{0}=1=9^{0}$, $D_{1}-D_{0}=10-1=9=9^{1}$, $D_{2}-D_{1}=91-10=81=9^{2}$, $D_{3}-D_{2}=820-91=729=9^{3}$, $D_{4}-D_{3}=7381-820=6561=9^{4}$.Thus, $D_{n}=D_{0} + 9^{1}+9^{2}+ . . . +9^{n}=\sum_{i=0}^{n}9^{i}=\frac{(1)(9^{n+1}-1)}{9-1}=\frac{9^{n+1}-1}{8}$.

 Thus, the desired sum is $\sum_{n=1}^{\infty}\frac{1}{8\frac{9^{n+1}-1}{8}+1}=\sum_{n=1}^{\infty}\frac{1}{9^{n+1}-1+1} = \sum_{n=1}^{\infty}\frac{1}{9^{n+1}}$

 This is an infinite \href{https://artofproblemsolving.com/wiki/index.php?title=Geometric\_series}{geometric series} with first term $\frac{1}{81}$ and common ratio $\frac{1}{9}$. Thus, the sum is $\frac{\frac{1}{81}}{1-\frac{1}{9}}=\frac{\frac{1}{81}}{\frac{8}{9}}=\frac{9}{(81)(8)}=\frac{1}{(9)(8)}=\frac{1}{72}$.

 Thus, $p + q = 1 + 72 = \boxed{073}$.

 





\textbf{Solution 2 (Engineer's Induction - not rigorous)}


 Manually, we can find \( D_1 = 10 \), \( D_2 = 91 \), \( D_3 = 820 \).   \[ \frac{1}{8D_1+1} = \frac{1}{81}, \quad \frac{1}{8D_2+1} = \frac{1}{729}, \quad \frac{1}{8D_3+1} = \frac{1}{6561}. \]   We notice that \( 81 = 9^2 \), \( 729 = 9^3 \), and \( 6561 = 9^4 \), and believe that this cannot be a coincidence.    Assume \( \frac{1}{8D_n+1} = \frac{1}{9^{n+1}} \). Then this problem asks us to find the sum of a geometric series with first term \( \frac{1}{81} \) and common ratio \( \frac{1}{9} \).   \[ \frac{\frac{1}{81}}{1 - \frac{1}{9}} = \frac{\frac{1}{81}}{\frac{8}{9}} = \frac{1}{72}. \]    Thus, the answer is \(\boxed{73}\).

 





\textbf{Solution 3 (Uses first half of Solution 1)}

From Solution 1, \( D_n = 10D_{n-1} - 9D_{n-2} \). From Solution 1, we also know \( D_2 = 91 \).  We can manually compute \( D_1 = 10 \). Iterating backwards, \(D_0 = 1 \).

 The characteristic polynomial of this recurrence is  
 \[x^2 - 10x + 9,\]  which has roots \( 1 \) and \( 9 \).  Therefore, the explicit formula for \( D_n \) is:  
 \[D_n = \alpha_1 + 9^n \alpha_2.\] 

 We can solve for \( \alpha_1 \) and \( \alpha_2 \) by setting \( n=0 \) and \( n=1 \) respectively. This gives the following equations:  
 \[\alpha_1 + \alpha_2 = 1,\] 
 \[\alpha_1 + 9\alpha_2 = 10.\] 

 Solving these equations, we find  
 \[\alpha_1 = -\frac{1}{8}, \quad \alpha_2 = \frac{9}{8}.\] 

 Substituting these back into the explicit formula, we find  
 \[D_n = \frac{9^{n+1} - 1}{8}.\] 

 Notice that  
 \[\frac{1}{8D_{n+1} + 1} = 9^{-(n+1)}.\] 

 The summation asks us to find the sum of a geometric series with ratio \( \frac{1}{9} \) and first term \( \frac{1}{81} \).  The sum of this series is  
 \[\frac{\frac{1}{81}}{1 - \frac{1}{9}} = \frac{\frac{1}{81}}{\frac{8}{9}} = \frac{1}{72}.\] 

 Thus, the answer is  $\boxed{73}$

 \textbf{Faster Alternate} Optionally, one can notice that the recursive relation is a first order non-homogenous relation in the form $a_n = ra_{n-1} + c$. So, to solve this (and all other first orders), we first need to find a point $L$ for which

 
 \[L = rL+c\]

 in which we have

 
 \[L = 9L+1\]
 \[L = -\frac{1}{8}\]

 Such $L$ is called the $\textbf{Equilibrium Point}$ of a recurrence relation. Its property satisfies:

 
 \[D_n - L = r(D_{n-1}-L)\]
 \[D_n + \frac{1}{8} = 9(D_{n-1}+\frac{1}{8})\]

 Let $b_n = D_n + \frac{1}{8}$. The result is a geometric sequence with first term $\frac{81}{8}$. The explicit rule for $b_n$ is just $b_n = \frac{81}{8} \cdot 9^{n-1}$. Conversion back to $D_n$ can be done quite easily.

 
 \[D_n + \frac{1}{8} = \frac{9^{n+1}}{8}\]
 \[D_n = \frac{9^{n+1} - 1}{8}\].

 The explicit rule has been reached and one can simply continue as follows.

 




	\item (\textit{AIME II 2011}) Nine delegates, three each from three different countries, randomly select chairs at a round table that seats nine people. Let the probability that each delegate sits next to at least one delegate from another country be $\frac{m}{n}$, where $m$ and $n$ are relatively prime positive integers. Find $m + n$.


 
 \textbf{Answer: }097



\textbf{Solution}

Use complementary probability and \href{https://artofproblemsolving.com/wiki/index.php?title=Principle\_of\_Inclusion-Exclusion}{Principle of Inclusion-Exclusion}. If we consider the delegates from each country to be indistinguishable and number the chairs, we have 
 \[\frac{9!}{(3!)^3} = \frac{9\cdot8\cdot7\cdot6\cdot5\cdot4}{6\cdot6} = 6\cdot8\cdot7\cdot5 = 30\cdot56\] total ways to seat the candidates. 

 Of these, there are $3 \times 9 \times \frac{6!}{(3!)^2}$ ways to have the candidates of at least some one country sit together. This comes to 
 \[\frac{27\cdot6\cdot5\cdot4}6 = 27\cdot 20.\]

 Among these there are $3 \times 9 \times 4$ ways for candidates from two countries to each sit together. This comes to $27\cdot 4.$

 Finally, there are $9 \times 2 = 18.$ ways for the candidates from all the countries to sit in three blocks (9 clockwise arrangements, and 9 counter-clockwise arrangements).

 So, by PIE, the total count of unwanted arrangements is $27\cdot 20 - 27\cdot 4 + 18 = 16\cdot27 + 18 = 18\cdot25.$ So the fraction 
 \[\frac mn = \frac{30\cdot 56 - 18\cdot 25}{30\cdot 56} = \frac{56 - 15}{56} = \frac{41}{56}.\] Thus $m + n = 56 + 41 = \fbox{097}.$

 


 \textbf{Setup Solution} We should give a specific setup for the table so to simplify. Call countries A, B, C. Someone from A (let me call him/her the "Master A") sits at the northernmost place for configurations. In my solution I choose to represent people by countries, because all the problem says about distinction among delegates is their country. Now, we can see that the total number of ways to arrange is nCr(8, 2) * nCr(6, 3) = 560 (8 seats choose 2 for the rest of the A country, 6 seats choose 3 for B).

 It is not difficult to see that the configurations we don't want are those which put three delegates in a row. Note also that despite my (arbitrary) setting A at the front, cases for A, B, C to be together are symmetric. Now, we use some \href{https://artofproblemsolving.com/wiki/index.php?title=Principle\_of\_Inclusion-Exclusion}{Principle of Inclusion-Exclusion}. For A three in a row: two cases. First. A, Master A, A: nCr(6, 3) for B seats = $20$. There's also $20$ for each of the "slant" A sets (i.e. A, A, Master A and Master A, A, A). =$60. *3=180$.For both A and B: Two cases again! A, Master A, A: we see 4 ways for B to have 3-in-a-row. Same for slant = $12. *3$ (AB, BC, or CA) = $36$.For all three: This time, we have 2 for each of the 3 aforementioned A situations = 6. Finally, 180-36+6=150. Those are the ones we don't want. The ones we do want are 410 of 560, so our answer is $\fbox{097}$.

 


 





\textbf{Solution 3}

We use complementary counting.

 We can order the $9$ people around a circle in $\frac{9!}{9} = 8!$ ways. Now we count when there is at least one delegate surrounded by people from only his/her country.

 Let the countries be $A,B,C$. Suppose that the group $XXX$ (for some country $X$) appears. To account for circular over counting we fix this group at the top. There are $6!$ ways to arrange the rest of the delegates and $3!$ ways to arrange inside the group. Since there are three countries this group can belong to, the total is $6!*3!*3$.

 But notice that when the group $XXX$ AND $YYY$ both appear is over counted. So, fix one group at the top. for the last country, we can insert $0,1,2,3$ people in between the two groups. Since we can choose two countries in $\binom{3}{2} = 3$ ways, the total is $3!*3!*3!*4*3$ ways.

 Unfortunately, there is over count in the over count. If each country has their delegates together, this case must be added back. Think of these three groups as a whole, and there are $3!/3 = 2$ ways of arranging. In each group there is $3!$ ways of arrangements. So, $3!*3!*3!*2$ is the total for this case.

 We now find the complementary probability total. This is 
 \[\frac{6!*3!*3 - 3!*3!*3!*4*3 + 3!*3!*3!*2}{8!} = \frac{15}{56}\] so the actual probability is $1-\frac{15}{56} = \frac{41}{56}$ for an answer of $\boxed{097}$.

 





\textbf{Solution 4 (PIE)}

$1680= \binom{9}{3} \binom{6}{3}$ is the total.

 $540 = 3* 9* \binom{6}{3}$ is the case where one country has three people in a row. (Three is for selection of country and nine is for rotation of the seats.)

 $108 = 3* 9* 4$ is the case where two countries has three people in a row. (Three and nine are for the same reasons above and four is for putting three people in a row for the remaining six seats.)

 $18$ is the case where three countries has three people in a row.

 
$\frac{1680-540+108-18}{1680}=\frac{41}{56}$ so the answer is $\boxed{097}$

 
By maxamc

 




	\item (\textit{AIME II 2011}) Point $P$ lies on the diagonal $AC$ of square $ABCD$ with $AP > CP$. Let $O_1$ and $O_2$ be the circumcenters of triangles $ABP$ and $CDP$, respectively. Given that $AB = 12$ and $\angle O_1PO_2 = 120 ^{\circ}$, then $AP = \sqrt{a} + \sqrt{b}$, where $a$ and $b$ are positive integers. Find $a + b$.


 
 \textbf{Answer: }096



\textbf{Solution 1}


                    \begin{figure}[H]
                        \centering
                        \includegraphics[width=0.5\textwidth]{Solution/AIME_II_2011_13_Solution_Figure_1.png}
                    \end{figure}
                    Note that $\angle BO_1P = 2\angle BAP = 90^\circ$. Since, $BO_1=PO_1$, this implies $\angle BPO_1 =45^\circ$. Similarly, $\angle  DPO_2= 45^\circ$. Thus, 
 \[\angle DPB = \angle O_1PO_2 + \angle BPO_1 -\angle DPO_2 = 120^\circ\] But, by SAS, $\triangle BAP \cong \triangle DAP$, so $\angle DPA = 60^\circ$, which implies $\angle ADP = 75^\circ$. Now we just apply law of sines:
 \[\frac{AP}{\sin 75^\circ} =\frac{12}{\sin 60^\circ} \implies AP = \sqrt{72}+\sqrt{24} \implies \boxed{096}\]Alternatively, you can note that $\triangle DAP$ is a 45-60-75 triangle, and solve from there.

 -sillybone

 





\textbf{Solution 2}

Denote the \href{https://artofproblemsolving.com/wiki/index.php?title=Midpoint}{midpoint} of $\overline{DC}$ be $E$ and the midpoint of $\overline{AB}$ be $F$. Because they are the circumcenters, both Os lie on the \href{https://artofproblemsolving.com/wiki/index.php?title=Perpendicular\_bisector}{perpendicular bisectors} of $AB$ and $CD$ and these bisectors go through $E$ and $F$.

 It is given that $\angle O_{1}PO_{2}=120^{\circ}$. Because $O_{1}P$ and $O_{1}B$ are \href{https://artofproblemsolving.com/wiki/index.php?title=Radius}{radii} of the same circle, the have the same length. This is also true of $O_{2}P$ and $O_{2}D$. Because $m\angle CAB=m\angle ACD=45^{\circ}$, $m\stackrel{\frown}{PD}=m\stackrel{\frown}{PB}=2(45^{\circ})=90^{\circ}$. Thus, $O_{1}PB$ and $O_{2}PD$ are isosceles right triangles. Using the given information above and symmetry, $m\angle DPB = 120^{\circ}$. Because ABP and ADP share one side, have one side with the same length, and one equal angle, they are congruent by SAS. This is also true for triangle CPB and CPD. Because angles APB and APD are equal and they sum to 120 degrees, they are each 60 degrees. Likewise, both angles CPB and CPD have measures of 120 degrees.

 Because the interior angles of a triangle add to 180 degrees, angle ABP has measure 75 degrees and angle PDC has measure 15 degrees. Subtracting, it is found that both angles $O_{1}BF$ and $O_{2}DE$ have measures of 30 degrees. Thus, both triangles $O_{1}BF$ and $O_{2}DE$ are 30-60-90 right triangles. Because F and E are the midpoints of AB and CD respectively, both FB and DE have lengths of 6. Thus, $DO_{2}=BO_{1}=4\sqrt{3}$. Because of 45-45-90 right triangles, $PB=PD=4\sqrt{6}$.

 Now, letting $x = AP$ and using \href{https://artofproblemsolving.com/wiki/index.php?title=Law\_of\_Cosines}{Law of Cosines} on $\triangle ABP$, we have

 
 \[96=144+x^{2}-24x\frac{\sqrt{2}}{2}\]
 \[0=x^{2}-12x\sqrt{2}+48\]

 Using the quadratic formula, we arrive at

 
 \[x = \sqrt{72} \pm \sqrt{24}\]

 Taking the positive root, $AP=\sqrt{72}+ \sqrt{24}$ and the answer is thus $\framebox[1.5\width]{096.}$

 NOTE: the positive root is taken because the negative root results in a degenerate triangle ABP.

 





\textbf{Solution 3}

This takes a slightly different route than Solution 1.

 Solution 1 proves that $\angle{DPB}=120^{\circ}$ and that $\overline{BP} = \overline{DP}$.Construct diagonal $\overline{BD}$ and using the two statements above it quickly becomes clear that $\angle{BDP} = \angle{DBP} = 30^{\circ}$ by isosceles triangle base angles.Let the midpoint of diagonal $\overline{AC}$ be $M$, and since the diagonals are perpendicular, both triangle $DMP$ and triangle $BMP$ are 30-60-90 right triangles.Since $\overline{AB} = 12$, $\overline{AC} = \overline{BD} = 12\sqrt{2}$ and $\overline{BM} = \overline{DM} = 6\sqrt{2}$.30-60-90 triangles' sides are in the ratio $1 : \sqrt{3} : 2$, so $\overline{MP} = \frac{6\sqrt{2}}{\sqrt{3}} = 2\sqrt {6}$.$\overline{AP} = \overline{MP} + \overline{BM} = 6\sqrt{2} + 2\sqrt{6} = \sqrt{72} + \sqrt{24}$.Hence, $72 + 24 = \framebox[1.5\width]{096}$.

 





\textbf{Solution 4}

Use vectors. In an $xy$ plane, let $(-s,0)$ be $A$, $(0,s)$ be $B$, $(s,0)$ be $C$, $(0,-s)$ be $D$, and $(p,0)$ be P, where $s=|AB|/\sqrt{2}=6\sqrt{2}$. It remains to find $p$.

 The line $y=-x$ is the \href{https://artofproblemsolving.com/wiki/index.php?title=Perpendicular\_bisector}{perpendicular bisector} of $AB$ and $CD$, so $O_1$ and $O_2$ lies on the line. Now compute the \href{https://artofproblemsolving.com/wiki/index.php?title=Perpendicular\_bisector}{perpendicular bisector} of $AP$. The center has coordinate $(\frac{p-s}{2},0)$, and the segment is part of the $x$-axis, so the perpendicular bisector has equation $x=\frac{p-s}{2}$. Since $O_1$ is the \href{https://artofproblemsolving.com/wiki/index.php?title=Circumcenter}{circumcenter} of triangle $ABP$, it lies on the perpendicular bisector of both $AB$ and $AP$, so
 \[O_1=(\frac{p-s}{2},-\frac{p-s}{2})\]Similarly, 
 \[O_2=(\frac{p+s}{2},-\frac{p+s}{2})\]The relation $\angle O_1PO_2=120^\circ$ can now be written using \href{https://artofproblemsolving.com/wiki/index.php?title=Vectors}{dot product} as
 \[\vec{PO_1}\cdot\vec{PO_2}=|\vec{PO_1}|\cdot|\vec{PO_2}|\cos 120^\circ=-\frac{1}{2}|\vec{PO_1}|\cdot|\vec{PO_2}|\]Computation of both sides yields
 \[\frac{p^2-s^2}{p^2+s^2}=-\frac{1}{2}\]Solve for $p$ gives $p=s/\sqrt{3}=2\sqrt{6}$, so $AP=s+p=6\sqrt{2}+2\sqrt{6}=\sqrt{72}+\sqrt{24}$. The answer is 72+24$\Rightarrow\boxed{096}$

 





\textbf{Solution 5}

Translate $\triangle{ABP}$ so that the image of $AB$ coincides $DC$. Let the image of $P$ be $P’$. 

 $\angle{DPC}=\angle{CPB}$ by symmetry, and $\angle{APB}=\angle{DP’C}$ because translation preserves angles. Thus $\angle{DP’C}+\angle{CPD}=\angle{CPB}+\angle{APB}=180^\circ$. Therefore, quadrilateral $CPDP’$ is cyclic. Thus the image of $O_1$ coincides with $O_2$. 

 $O_1P$ is parallel to $O_2P’$ so $\angle{P’O_2P}=\angle{O_1PO_2}=120^\circ$, so $\angle{PDP’}=60^\circ$ and $\angle{PDC}=15^\circ$, thus $\angle{ADP}=75^{\circ}$. 

 Let $M$ be the foot of the perpendicular from $D$ to $AC$. Then $\triangle{AMD}$ is a 45-45-90 triangle and $\triangle{DMP}$ is a 30-60-90 triangle. Thus

 $AM=6\sqrt{2}$ and $MP=\frac{6\sqrt{2}}{\sqrt{3}}$.

 This gives us $AP=AM+MP=\sqrt{72}+\sqrt{24}$, and the answer is $72+24=\boxed{096}.$

 





\textbf{Solution 6}

Reflect $O_1$ across $AP$ to $O_1'$. By symmetry $O_1’$ is the circumcenter of $\triangle{ADP}$

 $\angle{DO_1’P}$ = $2*\angle{DAP} = 90^\circ$, so $\angle{O_1’PD}=45^\circ$

 similarly $\angle{DO_2P}$ = $2*\angle{DCP} = 90^\circ$, so $\angle{O_2PD}=45^\circ$

 Therefore $\angle{O_1’PO_2}=90^\circ$, so that $\angle{O_1’PO_1} =120^\circ - 90^\circ = 30^\circ$

 By symmetry, $\angle{O_1'PA} = \angle{APO_1} = 0.5*\angle{O_1’PO_1} = 15^\circ$

 Therefore, since $O_1’$ is the circumcenter of $\triangle{ADP}$, $\angle{ADP}$ = $0.5*(180^\circ - 2*\angle{O_1'PA}) = 75^\circ$

 Therefore $\angle{APD} = 180^\circ - 45^\circ - 75^\circ = 60^\circ$

 Using sine rule in $\triangle{ADP}$, $AP = (12 * \sin 75^\circ) / \sin 60^\circ =\sqrt{72}+\sqrt{24}$, and the answer is $72+24=\boxed{096}.$

 By Kris17

 


 \href{https://mathproblemsolvingskills.wordpress.com/}{MathProblemSolvingSkills.com}

 


 





\textbf{Solution 7 (Coordinates)}

Why not use coordinates? After all, 45 degrees is rather friendly in terms of ordered-pair representation! We can set $A=(0, 12)$, $B=(12,12)$, $C=(12, 0)$, $D=(0, 0)$. Let this $P=(a, 12-a)$ for some $a$. 

 We also know that the circumcenter is the intersection of all perpendicular bisectors of sides, but two will suffice also due to this property. Therefore, we see that $O_{1}$ is the intersection of $x=6$ and, knowing the midpoint of $AP$ to be $(\frac{a}{2}, \frac{12-a}{2})$ and thus the equation to be $y=x+(12-a)$, we get $(6, 18-a)$. Likewise for $O_{2}$ it's $(6, 6-a)$. Now what do we see? $O_{1}P=O_{2}P$ (just look at the coordinates)! So both of those distances are $4\sqrt{3}$. Solving for $a$ we get it to be $6+2\sqrt{3}$, since $AP>CP$. Multiply by $\sqrt{2}$ because we are looking for $AP$ to get the answer of $\boxed{096}$.

 





\textbf{Solution 8 (Pure Angle Chasing)}

Let $\angle APD = \theta$. Then $\angle ADP = 180^{\circ}-45^{\circ}-\theta=135^{\circ}-\theta \implies \angle PDC=\theta-45^{\circ}$.Realize that because $O_1$ is a circumcenter, $\angle DO_1P=\angle DCP=45^{\circ} \implies \angle DPO_1=\frac{180^{\circ}-\angle DO_1P}{2}=45^{\circ}$. Then $\angle O_2PA=75^{\circ}-\theta \implies \angle AOP = 180^{\circ}-2\angle O_2PA = 2\theta+30^{\circ} \implies \angle ABP=\theta + 15^{\circ} \implies \angle PBC=75-\theta$.Now, because $P$ lies on diagonal $AC$, $\triangle PDC \cong \triangle PBC \implies \angle PDC = \angle PBC \implies \theta-45^{\circ}=75^{\circ}-\theta \implies \theta = 60^{\circ}$.To finish, we look at $\triangle ADP$. Drop a perpendicular from $D$ to $AP$ at $E$. Then $\triangle ADE$ is a $45-45-90$ and $\triangle PDE$ is a $30-60-90$. Therefore, $DE=EA=6\sqrt{2}, EP=2\sqrt{6}$, so $AP=AE+EP=6\sqrt{2}+2\sqrt{6}=\sqrt{72}+\sqrt{24} \implies \boxed{096}$. $\blacksquare$ ~msc

 





\textbf{Solution 9 (similar to solution 5)}

Both $O_1$ and $O_2$ lie on the perpendicular bisector of $AB$.

 \textbf{Claim:} $O_1O_2=12$ and $O_1P=O_2P$.

 \textit{Proof.} Translate $O_1$ and $P$ $12$ units down, and let their images be $O_1'$ and $P'$, respectively. Note that $\triangle ABP\cong\triangle DCP'$. Additionally, 
 \[\angle CP'D = \angle BPA = 180^{\circ} - \angle BPC = 180^{\circ} - \angle CPD,\] so $CPDP'$ is cyclic. This means $O_1'$ and $O_2$ coincide, so $O_1O_2=12$. This also means the circumradii of both triangles are equal, so $O_1P=O_2P$. $\blacksquare$.

 Let the perpendicular from $P$ intersect $O_1O_2$ at $X$ and $AD$ at $Y$. Since $\triangle O_1XP$ is 30-60-90, $XP=\frac{6}{\sqrt{3}} = 2\sqrt3$. Since $YX=6$, $PY=6+2\sqrt3$, so $AP=6\sqrt2+2\sqrt6 = \sqrt{72}+\sqrt{24} \implies\boxed{96}$.

 





\textbf{Solution 10}


                    \begin{figure}[H]
                        \centering
                        \includegraphics[width=0.5\textwidth]{Solution/AIME_II_2011_13_Solution_Figure_2.png}
                    \end{figure}
                    
 \[BP = DP,    \angle PAB =\angle PCD  = 45^\circ \implies  O_{1}P =O_{2}P\]  by the Law of Sines.

 
 \[AB= CD, O_{1}A = O_{2}D = O_{1}B = O_{2}C \implies \triangle AO_{1}B = \triangle DO_{2}C.\] They have translational symmetry. The translation vector is 
 \[\vec {AD} \implies  O_{1}O_{2} = AD.\] 

 
 \[\triangle PO_{1}O_{2} = \triangle O_{1}AB \implies \angle AO_{1}B = 120^\circ \implies\]
 \[\angle APB = 60^\circ \implies  \angle ABP = 180^\circ – 60^\circ – 45^\circ = 75^\circ.\]By the Law of Sines 
 \[AP = \frac {AB \cdot \sin 75^\circ}{\sin 60^\circ} = \frac {AB \cdot \sin (30^\circ + 45^\circ)} {\sin 60^\circ}\]
 \[AP  = AB \cdot (\sin 45^\circ + \cos 45^\circ \cdot \tan 30^\circ),\]
 \[AP = \frac {AB}{\sqrt{2}} (1 + \frac {1}{\sqrt{3}}) = 6 \sqrt {2} + 2 \sqrt {6} \implies \boxed{\textbf{096}.}\]\textbf{vladimir.shelomovskii@gmail.com, vvsss}

 




	\item (\textit{AIME II 2011}) There are $N$ permutations $(a_1, a_2, \dots, a_{30})$ of $1, 2, \dots, 30$ such that for $m \in \{2,3,5\}$, $m$ divides $a_{n+m} - a_n$ for all integers $n$ with $1 \le n < n+m \le 30$. Find the remainder when $N$ is divided by 1000.


 
 \textbf{Answer: }440



\textbf{Solution 1}

Be wary of "position" versus "number" in the solution!

 Each POSITION in the 30-position permutation is uniquely defined by an ordered triple $(i, j, k)$. The $n$th position is defined by this ordered triple where $i$ is $n \mod 2$, $j$ is $n \mod 3$, and $k$ is $n \mod 5$. There are 2 choices for $i$, 3 for $j$, and 5 for $k$, yielding $2 \cdot 3 \cdot 5=30$ possible triples. Because the least common multiple of 2, 3, and 5 is 30, none of these triples are repeated and all are used. By the conditions of the problem, if $i$ is the same in two different triples, then the two numbers in these positions must be equivalent mod 2. If $j$ is the same, then the two numbers must be equivalent $\mod 3$, and if $k$ is the same, the two numbers must be equivalent mod 5. Take care to note that that doesn't mean that the number 1 has to have $1 \mod 2$! It's that the POSITION which NUMBER 1 occupies has $1 \mod 2$!

 The ordered triple (or position) in which 1 can be placed has 2 options for i, 3 for j, and 5 for k, resulting in 30 different positions of placement.

 The ordered triple where 2 can be placed in is somewhat constrained by the placement of 1. Because 1 is not equivalent to 2 in terms of mod 2, 3, or 5, the i, j, and k in their ordered triples must be different. Thus, for the number 2, there are (2-1) choices for i, (3-1) choices for j, and (5-1) choices for k. Thus, there are 1*2*4=8 possible placements for the number two once the number one is placed.

 Because 3 is equivalent to 1 mod 2, it must have the same i as the ordered triple of 1. Because 3 is not equivalent to 1 or 2 in terms of mod 3 or 5, it must have different j and k values. Thus, there is 1 choice for i, (2-1) choices for j, and (4-1) choices for k, for a total of $1\cdot 1 \cdot 3=3$ choices for the placement of 3.

 As above, 4 is even, so it must have the same value of i as 2. It is also 1 mod 3, so it must have the same j value of 1. 4 is not equivalent to 1, 2, or 3 mod 5, so it must have a different k value than that of 1, 2, and 3. Thus, there is 1 choice for i, 1 choice for j, and (3-1) choices for k, yielding a total of $1 \cdot 1 \cdot 2=2$ possible placements for 4.

 5 is odd and is equivalent to 2 mod 3, so it must have the same i value as 1 and the same j value of 2. 5 is not equivalent to 1, 2, 3, or 4 mod 5, so it must have a different k value from 1, 2, 3, and 4. However, 4 different values of k are held by these four numbers, so 5 must hold the one remaining value. Thus, only one possible triple is found for the placement of 5.

 All numbers from 6 to 30 are also fixed in this manner. All values of i, j, and k have been used, so every one of these numbers will have a unique triple for placement, as above with the number five. Thus, after 1, 2, 3, and 4 have been placed, the rest of the permutation is fixed. 

 Thus, $N = 30 \cdot 8 \cdot 3 \cdot 2=30 \cdot 48=1440$. Thus, the remainder when $N$ is divided by $1000$ is $\boxed{440}.$

 





\textbf{Solution 2}

We observe that the condition on the permutations means that two numbers with indices congruent $\mod m$ are themselves congruent $\mod m$ for $m \in \{ 2,3,5\}.$ Furthermore, suppose that $a_n \equiv k \mod m.$ Then, there are $30/m$ indices congruent to $n \mod m,$ and $30/m$ numbers congruent to $k \mod m,$ because 2, 3, and 5 are all factors of 30. Therefore, since every index congruent to $n$ must contain a number congruent to $k,$ and no number can appear twice in the permutation, only the indices congruent to $n$ contain numbers congruent to $k.$ In other words, $a_i \equiv a_j \mod m \iff i \equiv j \mod m.$ But it is not necessary that $\textcolor{red}{(a_i\equiv i)\cup (a_j\equiv j)\mod m}$. In fact, if that were the case, there would only be one way to assign the indices, since $2,3,5$ are relatively prime to each other and $30=\text{lcm}(2,3,5)$: $\{a_1,a_2,\dots a_{30}\}\in\{1,2,\dots 30\}\text{ }respectively$.

 This tells us that in a valid permutation, the congruence classes $\mod m$ are simply swapped around, and if the set $S$ is a congruence class $\mod m$ for $m =$ 2, 3, or 5, the set $\{ a_i \vert i \in S \}$ is still a congruence class $\mod m.$ Clearly, each valid permutation of the numbers 1 through 30 corresponds to exactly one permutation of the congruence classes modulo 2, 3, and 5. Also, if we choose some permutations of the congruence classes modulo 2, 3, and 5, they correspond to exactly one valid permutation of the numbers 1 through 30. This can be shown as follows: First of all, the choice of permutations of the congruence classes gives us every number in the permutation modulo 2, 3, and 5, so by the Chinese Remainder Theorem, it gives us every number $\mod 2\cdot 3\cdot 5 = 30.$ Because the numbers must be between 1 and 30 inclusive, it thus uniquely determines what number goes in each index. Furthermore, two different indices cannot contain the same number. We will prove this by contradiction, calling the indices $a_i$ and $a_j$ for $i \neq j.$ If $a_i=a_j,$ then they must have the same residues modulo 2, 3, and 5, and so $i \equiv j$ modulo 2, 3, and 5. Again using the Chinese Remainder Theorem, we conclude that $i \equiv j \mod 30,$ so because $i$ and $j$ are both between 1 and 30 inclusive, $i = j,$ giving us a contradiction. Therefore, every choice of permutations of the congruence classes modulo 2, 3, and 5 corresponds to exactly one valid permutation of the numbers 1 through 30.

 In other words, each set of assignment from $a_j\rightarrow j\mod (2,3,5)$ determines a unique string of $30$ numbers. For example:

 $\left[(0,1)\rightarrow (1,0)\right]\cap\left[(0,1,2)\rightarrow (0,2,1)\right]\cap\left[(0,1,2,3,4)\rightarrow (4,2,3,0,1)\right]$:
 \[\begin{array}{|c||c|c|c|c|c|c|c|c|c|c|c|c|c|c|c|c|c|c|c|c|c|c|c|c|c|c|c|c|c|c|} \hline 2&1&0&1&0&1&0&1&0&1&0&1&0&1&0&1&0&1&0&1&0&1&0&1&0&1&0&1&0&1&0\\ \hline 3&0&2&1&0&2&1&0&2&1&0&2&1&0&2&1&0&2&1&0&2&1&0&2&1&0&2&1&0&2&1\\ \hline 5&4&2&3&0&1&4&2&3&0&1&4&2&3&0&1&4&2&3&0&1&4&2&3&0&1&4&2&3&0&1\\ \hline\hline 30&9&2&13&0&11&4&27&8&25&6&29&22&3&20&1&24&17&28&15&26&19&12&23&10&21&14&7&18&5&16\\ \hline \end{array}\]

 We have now established a bijection between valid permutations of the numbers 1 through 30 and permutations of the congruence classes modulo 2, 3, and 5, so $N$ is equal to the number of permutations of congruence classes. There are always $m$ congruence classes $\mod m,$ so $N = 2! \cdot 3! \cdot 5! = 2 \cdot 6 \cdot 120 = 1440 \equiv \framebox[1.3\width]{440} \mod 1000.$

 





\textbf{Solution 3 (2-sec solve)}

Note that $30=2\cdot 3\cdot 5$. Since $\gcd(2, 3, 5)=1$, by CRT, for each value $k=0\ldots 29$ modulo $30$ there exists a unique ordered triple of values $(a, b, c)$ such that $k\equiv a\pmod{2}$, $k\equiv b\pmod{3}$, and $k\equiv c\pmod{5}$. Therefore, we can independently assign the residues modulo $2, 3, 5$, so $N=2!\cdot 3!\cdot 5!=1440$, and the answer is $\boxed{440}$.

 -vockey

 





\textbf{Solution 4( Better explanation of Solution 3)}

First let's look at the situation when $m$ is equal to 2. It isn't too difficult to see the given conditions are satisfied iff the sequences  $S_1 = a_1,a_3,a_5,a_7...$ and $S_2 =a_2,a_4,a_6...$ each are assigned either $\equiv 0 \pmod 2$ or $\equiv 1 \pmod 2$. Another way to say this is each element in the sequence $a_1,a_3,a_5,a_7...$ would be the same mod 2, and similarly for the other sequence. There are 2! = 2 ways to assign the mods to the sequences.

 Now when $m$ is equal to 3, the sequences are $T_1 = a_1, a_4,a_7..$, $T_2 = a_2,a_5,a_8,a_11...$, and $T_3 = a_3,a_6,a_9...$. Again, for each sequence, all of its elements are congruent either $0$,$1$, or $2$ mod $3$. There are $3! = 6$ ways to assign the mods to the sequences.

 Finally do the same thing for $m = 5$. There are $5!$ ways. In total there are $2 * 6*120 = 1440$ and the answer is $\boxed{440}$.

 Why does this method work? Its due to CRT.

 -MathLegend27

 Note: the explanation is not complete without mentioning that 2*3*5 = 30, therefore CRT guarantees there is only 1 permutation for every combination of mod selections. If the problem asked for permutations of ${1,2,...,60}$, the answer would have been $2!3!5!2!2!2!$. 

 (I think for 60 u mulitply by $2^{30}$ not $(2!)^3$)

 -Mathdummy

 





\textbf{Solution 5(slightly more rigor than needed)}

Note, for the purposes of this problem, it is easier to have our residue classes mod \( m \) to be \( \{1, 2, \dots, m\} \) instead of \( \{0, 1, 2, \dots, m-1\} \).

 We will start by finding a general insight for \( m \) given that \( m \mid 30 \).  
 \[m \mid (a_{i + m} - a_i) \implies a_{i + m} \equiv a_i \pmod{m}\]  Notice that for every \( p \) between \( 1 \) and \( m \) inclusive, we have  
 \[a_p \equiv a_{p+m} \equiv a_{p+2m} \equiv a_{p+3m} ... \pmod{m}.\]  Notice that \( \{p, p+m, p+2m, \dots, p+km\} \) all fall under the residue class \( p \pmod{m} \).  Let this list of numbers with indices that are \( p \pmod{m} \) all be congruent to \( \sigma_m(p) \pmod{m} \).  Notice that \( \sigma_m \) maps numbers from \( \{1, 2, \dots, m\} \) to \( \{1, 2, \dots, m\} \).

 $\textbf{Claim: \( \sigma_m \) is bijective.}$Since \( \sigma_m \) is a valid function mapping two sets of the same size, we simply prove injectivity.  

 Assume for the sake of contradiction, for \( p \neq q \), let \( \sigma_m(p) \equiv \sigma_m(q) \pmod{m} \), and let them be congruent to some constant \( x \pmod{m} \).  Since \( 1 \leq p, q \leq m \) and \( p \neq q \), this implies \( p \) and \( q \) are not congruent mod \( m \).  
 \[\{p, p+m, p+2m, \dots \} \cap \{q, q+m, q+2m, \dots\} = \emptyset,\]  since these sets are distinct mod \( m \).  

 Therefore, \( \sigma_m \) maps \( \frac{2 \cdot 30}{m} \) distinct values to some \( x \pmod{30} \).  However, this implies that there are \( \frac{2 \cdot 30}{m} \) distinct values between \( 1 \) and \( 30 \) falling under the same residue class, which is obviously false.  Thus, we have proven the claim.  

 Since \( \sigma_m \) is injective and maps between two sets of identical size, it is bijective.  Therefore, there are \( m! \) ways to create such a mapping called \( \sigma_m \).  

 For this problem, there are \( 2! \cdot 3! \cdot 5! \) ways to find \( \sigma_2 \), \( \sigma_3 \), and \( \sigma_5 \).  (Notice \( \sigma_k \) lies in the symmetric group \( S_k \).)

 $\textbf{Claim: Finding \( \sigma_2 \), \( \sigma_3 \), and \( \sigma_5 \) completely determines \( a_i \).}$ If we know \( \sigma_2(i) \), \( \sigma_3(i) \), and \( \sigma_5(i) \), we know \( i \pmod{2} \), \( i \pmod{3} \), and \( i \pmod{5} \).  By the Chinese Remainder Theorem (CRT), we know \( i \pmod{30} \), and since \( 1 \leq i \leq 30 \), this completely determines \( i \).  

 Therefore, \( 2! \cdot 3! \cdot 5! = 2 \cdot 6 \cdot 120 = 1440 \).  

 Therefore, the answer is $\boxed{440}$

 \href{https://mathproblemsolvingskills.wordpress.com/}{MathProblemSolvingSkills.com}

 


 




	\item (\textit{AIME II 2011}) Let $P(x) = x^2 - 3x - 9$. A real number $x$ is chosen at random from the interval $5 \le x \le 15$. The probability that $\left\lfloor\sqrt{P(x)}\right\rfloor = \sqrt{P(\lfloor x \rfloor)}$ is equal to $\frac{\sqrt{a} + \sqrt{b} + \sqrt{c} - d}{e}$ , where $a$, $b$, $c$, $d$, and $e$ are positive integers. Find $a + b + c + d + e$.


 
 \textbf{Answer: }850



\textbf{Solution 1}

Table of values of $P(x)$:

 
 \begin{align*} P(5) &= 1 \\ P(6) &= 9 \\ P(7) &= 19 \\ P(8) &= 31 \\ P(9) &= 45 \\ P(10) &= 61 \\ P(11) &= 79 \\ P(12) &= 99 \\ P(13) &= 121 \\ P(14) &= 145 \\ P(15) &= 171 \\ \end{align*}


 In order for $\lfloor \sqrt{P(x)} \rfloor = \sqrt{P(\lfloor x \rfloor)}$ to hold, $\sqrt{P(\lfloor x \rfloor)}$ must be an integer and hence $P(\lfloor x \rfloor)$ must be a perfect square. This limits $x$ to $5 \le x < 6$ or $6 \le x < 7$ or $13 \le x < 14$ since, from the table above, those are the only values of $x$ for which $P(\lfloor x \rfloor)$ is a perfect square. However, in order for $\sqrt{P(x)}$ to be rounded down to $P(\lfloor x \rfloor)$, $P(x)$ must be less than the next perfect square after $P(\lfloor x \rfloor)$ (for the said intervals). Now, we consider the three cases:

 
Case $5 \le x < 6$:

 $P(x)$ must be less than the first perfect square after $1$, which is $4$, \textit{i.e.}:

 $1 \le P(x) < 4$ (because $\lfloor \sqrt{P(x)} \rfloor = 1$ implies $1 \le \sqrt{P(x)} < 2$)

 Since $P(x)$ is increasing for $x \ge 5$, we just need to find the value $v \ge 5$ where $P(v) = 4$, which will give us the working range $5 \le x < v$.

 
 \begin{align*} v^2 - 3v - 9 &= 4 \\ v &= \frac{3 + \sqrt{61}}{2} \end{align*}


 So in this case, the only values that will work are $5 \le x < \frac{3 + \sqrt{61}}{2}$.

 Case $6 \le x < 7$:

 $P(x)$ must be less than the first perfect square after $9$, which is $16$.

 
 \begin{align*} v^2 - 3v - 9 &= 16 \\ v &= \frac{3 + \sqrt{109}}{2} \end{align*}


 So in this case, the only values that will work are $6 \le x < \frac{3 + \sqrt{109}}{2}$.

 Case $13 \le x < 14$:

 $P(x)$ must be less than the first perfect square after $121$, which is $144$.

 
 \begin{align*} v^2 - 3v - 9 &= 144 \\ v &= \frac{3 + \sqrt{621}}{2} \end{align*}


 So in this case, the only values that will work are $13 \le x < \frac{3 + \sqrt{621}}{2}$.

 Now, we find the length of the working intervals and divide it by the length of the total interval, $15 - 5 = 10$:

 
 \begin{align*} \frac{\left( \frac{3 + \sqrt{61}}{2} - 5 \right) + \left( \frac{3 + \sqrt{109}}{2} - 6 \right) + \left( \frac{3 + \sqrt{621}}{2} - 13 \right)}{10} \\ &= \frac{\sqrt{61} + \sqrt{109} + \sqrt{621} - 39}{20} \end{align*}


 Thus, the answer is $61 + 109 + 621 + 39 + 20 = \fbox{850}$.

 





\textbf{Solution 2}

Make the substitution $y=2x-3$, so $P(x)=\frac{y^2-45}{4}.$ We're looking for solutions to 
 \[\bigg\lfloor{\sqrt{\frac{y^2-45}{4}}\bigg\rfloor}=\sqrt{\frac{\lfloor{y\rfloor}^2-45}{4}}\]with the new bounds $y\in{[7,27]}$. Since the left side is an integer, it must be that $\frac{\lfloor{y\rfloor}^2-45}{4}$ is a perfect square. For simplicity, write $\lfloor{y\rfloor}=a$ and 
 \[a^2-45=4b^2\implies{(a-2b)(a+2b)=45}.\]Since $a-2b<a+2b$, it must be that $(a-2b,a+2b)=(1,45),(3,15),(5,9)$, which gives solutions $(23,11),(9,3),(7,1)$, respectively. But this gives us three cases to check:

 Case 1: $\bigg\lfloor{\sqrt{\frac{y^2-45}{4}}\bigg\rfloor}=11$.

 In this case, we have 
 \[11\leq{\sqrt{\frac{y^2-45}{4}}}<12\implies{y\in{[23,\sqrt{621})}}.\]Case 2: $\bigg\lfloor{\sqrt{\frac{y^2-45}{4}}\bigg\rfloor}=3$.

 In this case, we have 
 \[3\leq{\sqrt{\frac{y^2-45}{4}}}<4\implies{y\in{[9,\sqrt{109})}}.\]Case 3: $\bigg\lfloor{\sqrt{\frac{y^2-45}{4}}\bigg\rfloor}=1$

 In this case, we have 
 \[1\leq{\sqrt{\frac{y^2-45}{4}}}<2\implies{y\in{[7,\sqrt{61})}}.\]To finish, the total length of the interval from which we choose $y$ is $27-7=20$. The total length of the success intervals is 
 \[(\sqrt{61}-7)+(\sqrt{109}-9)+(\sqrt{621}-23)=\sqrt{61}+\sqrt{109}+\sqrt{621}-39,\]which means the probability is 
 \[\frac{\sqrt{61}+\sqrt{109}+\sqrt{621}-39}{20}.\]The requested sum is $\boxed{850}$.

 





\textbf{Solution 3 (Graphing)}

It's definitely possible to conceptualize this problem visually, if that's your thing, since it is a geometric probability problem. Let $A = \left\lfloor\sqrt{P(x)}\right\rfloor$ and $B = \sqrt{P(\left\lfloor x \right\rfloor)}$. The graph of $A$ and $B$ will look like this, with $A$ having only integral y-values and $B$ having only integral x-values:

 
                    \begin{figure}[H]
                        \centering
                        \includegraphics[width=0.5\textwidth]{Solution/AIME_II_2011_15_Solution_Figure_1.png}
                    \end{figure}
                    

 As both $A$ and $B$ consist of a bunch of line segments, the probability that $A = B$ is the "length" of the overlap between the segments of $A$ and $B$ divided by the total length of the segments of $B$. 

 Looking at the graph, we see that $A$ and $B$ will overlap only when $B$ is an integer. Specifically, each region of overlap will begin when $\sqrt{P(x)}\ = k,  5 \le x \le 15$ has solutions for integral $k$ in the range of $A$, which consists of the integers $1-13$, and end when $A$ jumps up to its next y-value.

 Using the quadratic formula, we see that the start point of each of these overlapping segments will be the integral values of $\frac{3 + \sqrt{45 + 4k^2}}{2}$ for $k$ in the specified range, meaning $45 + 4k^2$ must be a perfect square. Plugging in all the possible values of $k$, we get $k = 1, 3, 11$, corresponding to start points of $x = 5, 6, 13$. As already stated, the endpoints will occur when $A$ jumps up to the next integer $k+1$ at each of these segments, at which point the x-value will be $\frac{3 + \sqrt{45 + 4(k+1)^2}}{2}$. On the graph, the overlapping segments of $A$ and $B$ would be represented by the highlighted green segments below:

 
                    \begin{figure}[H]
                        \centering
                        \includegraphics[width=0.5\textwidth]{Solution/AIME_II_2011_15_Solution_Figure_2.png}
                    \end{figure}
                    

 
Taking the difference between this second x-value and the start point for each of our start points $x = 5, 6, 13$ and summing them will give us the total length of these green segments. We can then divide this value by ten (the total length of the segments of $B$) to give us the probability of overlap between $A$ and $B$.

 Doing so gives us:

 $\frac{\left( \frac{3 + \sqrt{61}}{2} - 5 \right) + \left( \frac{3 + \sqrt{109}}{2} - 6 \right) + \left( \frac{3 + \sqrt{621}}{2} - 13 \right)}{10} = \frac{\sqrt{61} + \sqrt{109} + \sqrt{621} - 39}{20}$

 $\implies{61 + 109 + 621 + 39 + 20 = \fbox{850}}$.

 





\textbf{Solution 4}

Note that all the "bounds" have to be less than the number+1, otherwise it wouldn't fit the answer format. Therefore, the answer is $\frac{3*3+\sqrt{9+4(4+9)}-10+\sqrt{9+4(16+9)}-12+\sqrt{9+4(144+9)}}{20} \implies \boxed{850}$

 \href{https://mathproblemsolvingskills.wordpress.com/}{MathProblemSolvingSkills.com}

 




\end{enumerate}